\chapimg{fig/sh06.jpg}
\IfFileExists{extra/lead06.tex}{\chapterlead{\subhd{After this chapter you can} classify irrigation methods and compute their efficiencies and uniformities; size a pipe by the Hazen--Williams and Darcy--Weisbach equations and find the pump head and power; use the continuity and Bernoulli equations; prepare and check weather data; estimate missing radiation; and compute vegetation indices from satellite bands.}
}{}
\chapter{Irrigation Design Mathematics, Weather and Field Data}\label{chap:06}
\markboth{Irrigation Design Mathematics, Weather and Field Data}{Irrigation Design Mathematics, Weather and Field Data}
\banner{Overview of the subject}
\IfFileExists{extra/bg06.tex}{\section{Introduction: designing a system that delivers the water}
The first five chapters answered how much water a soil holds, how it moves, when to irrigate and how much a crop uses. A farmer has still to get that water to the plant: from a well or a canal, through pipes or channels, into the soil, with what pressure and at what cost. This is \kw{irrigation design}. It rests on a few ideas from hydraulics (continuity, energy, friction), on a clear idea of efficiency (how much of the water taken from the source reaches the roots), and on data about the weather and the field that set the demand. This chapter collects these ideas; the later chapters apply them to surface, drip and sprinkler systems and to drainage.

\section{Classification of irrigation methods}
Irrigation methods divide into three families by the way the water reaches the soil. In \kw{surface irrigation} the water flows over the field by gravity: \emph{basin} irrigation floods a levelled area, \emph{border} irrigation uses long strips with a free downstream end, and \emph{furrow} irrigation runs the water in small channels between the crop rows. It needs no pressure, but the application is non-uniform unless the field is well levelled and the inflow well chosen; efficiencies are 50--70 per cent. In \kw{sprinkler irrigation} water is sprayed through nozzles under pressure: \emph{set} systems (solid set, hand-move, side-roll), \emph{moving} systems (centre pivot, lateral move, travelling gun) and \emph{rotating} or \emph{impact} sprinklers differ in cost and labour; efficiencies are 70--85 per cent. In \kw{micro-irrigation} (drip, micro-sprinkler, subsurface drip) small flows are given to each plant through emitters at low pressure; efficiencies are 85--95 per cent and the wetted volume is a fraction of the soil, which suits row crops, trees and vines and needs filtration to prevent clogging.

\section{Efficiency and uniformity}
Not all water taken from the source benefits the crop. Four efficiencies are used. The \kw{conveyance efficiency} $E_c$ is the ratio of the water that reaches the field to the water taken at the source; the \kw{application efficiency} $E_a=V_{stored}/V_{applied}$ is the share of the applied water stored in the root zone; the \kw{storage efficiency} $E_s=V_{stored}/V_{needed}$ is the share of the needed depth that has been supplied; and the \kw{distribution uniformity} $DU_{lq}$ is the mean depth of the lowest quarter divided by the mean depth. The overall field efficiency is $E=E_cE_a$. The \kw{Christiansen coefficient} $CU=1-\sum|z_i-\bar z|/(n\bar z)$ is the other standard measure of uniformity, and, for emitters, the \kw{emission uniformity} $EU=q_{min}/\bar q$.

\exbox{Uniformity from catch cans}{Twelve catch cans in a sprinkler field collect (mm): 12, 14, 15, 15, 16, 17, 17, 18, 19, 22, 10, 13. Find the mean, Christiansen's coefficient, and the distribution uniformity of the lowest quarter.}\sol Sum $=188$, mean $\bar z=15.67$ mm. Absolute deviations sum to $30.00$; $CU=1-30.00/(12\times15.67)=0.840$ \ans Lowest quarter: the three smallest values $10, 12, 13$, mean $11.67$; $DU_{lq}=11.67/15.67=0.745$ \finans

\section{Hydraulics: continuity, energy and friction}
Pressurised systems obey three relations. \kw{Continuity}: the flow $Q=vA$ is the same at every section of a pipe without branches, so a pipe of diameter $D$ carries the velocity $v=4Q/(\pi D^2)$. \kw{Energy} (Bernoulli): the total head $H=z+p/(\rho g)+v^2/(2g)$ falls along a pipe by the friction head loss $h_f$ and the local losses $h_m=K_mv^2/(2g)$ of fittings. \kw{Friction}: for water in a smooth plastic pipe the empirical Hazen--Williams equation in SI units is

\begin{equation}h_f=10.67\,\frac{L\,Q^{1.852}}{C^{1.852}\,D^{4.871}},\label{eq:hw6}\end{equation}


with $Q$ in m$^3$\,s$^{-1}$, $D$ in m, $L$ in m and $C$ the roughness coefficient (140 for new PVC and polyethylene pipe, 150 for the smooth laterals; lower for old steel). The Darcy--Weisbach form $h_f=f\,(L/D)\,v^2/(2g)$ with the friction factor $f$ from the Moody diagram is more general. A pipe is sized so that the velocity lies between 0.3 m\,s$^{-1}$ (to keep sediment moving) and 1.5 m\,s$^{-1}$ (to limit water hammer), and so that the head loss in the pipe and the cost of the pipe together are the least. The power of a pump that delivers $Q$ at a total dynamic head $H$ with an efficiency $\eta$ is $P=\rho g Q H/\eta$.

\exbox{Mainline friction and power}{A PVC mainline of $L=600$ m and $D=200$ mm ($C=140$) carries $Q=45$ L\,s$^{-1}$. Find the velocity, the head loss, and the extra pump power needed to overcome it ($\eta=0.65$).}\sol \givenline{$L=600$ m ; $D=200$ mm ; $C=140$ ; $Q=45$ L\,s$^{-1}$ ; $\eta=0.65$}\ans $Q=0.045$ m$^3$\,s$^{-1}$; $v=4\times0.045/(\pi\times0.2^2)=1.43$ m\,s$^{-1}$ (within 0.3--1.5). Head loss: $10.67\times600\times0.045^{1.852}/(140^{1.852}\times0.2^{4.871})=5.52$ m \ans Power $=\rho gQh_f/\eta=1000\times9.81\times0.045\times5.52/0.65=3750$ W $=3.75$ kW \finans

\section{Weather data for a design}
Every design and every model needs the weather: rainfall, maximum and minimum temperature, solar radiation (or sunshine hours), humidity (or dew point) and wind speed, daily, for many years. The data may come from a weather station near the field, from the national service, from a gridded reanalysis or from a satellite product. Three checks are needed before use: \emph{completeness} (gaps are filled by interpolation from neighbouring days or stations), \emph{plausibility} (negative rain, $T_{min}>T_{max}$, radiation above the extraterrestrial value are errors) and \emph{homogeneity} (a change of instrument shows as a step in the mean). When radiation is missing it may be estimated from the temperature range by Hargreaves' relation $R_s=k_{Rs}\sqrt{T_{max}-T_{min}}\,R_a$ with $k_{Rs}\approx0.16$ (interior) to 0.19 (coastal) and $R_a$ the extraterrestrial radiation. For a design the \emph{dependable} (e.g. 80 per cent probability) rainfall and the \emph{peak} $ET_0$ of the design period are used, not the means.

\exbox{Radiation from the temperature range}{A station lacks radiation. For $T_{max}=31$, $T_{min}=17$ $^\circ$C and an extraterrestrial radiation $R_a=38$ MJ\,m$^{-2}$\,d$^{-1}$ (mid-latitude summer day), estimate $R_s$ with $k_{Rs}=0.17$ and the net shortwave radiation for an albedo 0.23.}\sol \givenline{$T_{max}=31$ ; $T_{min}=17$ $^\circ$C ; $R_a=38$ MJ\,m$^{-2}$\,d$^{-1}$ ; $k_{Rs}=0.17$}\ans $R_s=0.17\sqrt{31-17}\times38=0.17\times3.742\times38=24.2$ MJ\,m$^{-2}$\,d$^{-1}$ \ans $R_{ns}=(1-0.23)R_s=18.6$ MJ\,m$^{-2}$\,d$^{-1}$ \ans The clear-sky bound is $0.75R_a=28.5$, so the estimate is plausible. \finans

\section{Field data from satellites: vegetation indices}
A satellite measures the reflectance of the surface in bands of the spectrum. Healthy green vegetation absorbs red light (chlorophyll) and reflects near-infrared strongly (leaf structure); bare soil reflects both moderately; water absorbs both. The \kw{normalised difference vegetation index} is

\begin{equation}\mathrm{NDVI}=\frac{\rho_{NIR}-\rho_{red}}{\rho_{NIR}+\rho_{red}},\label{eq:ndvi6}\end{equation}
\par\noindent where NDVI is the normalised difference vegetation index, $\rho_{NIR}$ the reflectance in the near-infrared band and $\rho_{red}$ the reflectance in the red band (both 0 to 1).



from $-1$ (water) through 0 (bare soil) to about 0.9 (dense canopy). Other indices use other bands: the \emph{normalised difference water index} with a shortwave band follows the water in the canopy, and the \emph{soil-adjusted} index corrects for soil brightness. The fraction of the ground covered by the canopy is estimated from NDVI by $f_c=[(\mathrm{NDVI}-\mathrm{NDVI}_s)/(\mathrm{NDVI}_v-\mathrm{NDVI}_s)]^2$, and the basal crop coefficient from $f_c$; this is how an image can drive a schedule.

\exbox{NDVI and cover}{A pixel has reflectances $\rho_{red}=0.08$ and $\rho_{NIR}=0.42$. Find the NDVI and the cover fraction if $\mathrm{NDVI}_s=0.15$ (bare) and $\mathrm{NDVI}_v=0.85$ (full cover).}\sol \givenline{$\rho_{red}=0.08$ ; $\rho_{NIR}=0.42$ ; $\mathrm{NDVI}_s=0.15$ ; $\mathrm{NDVI}_v=0.85$}\ans $\mathrm{NDVI}=(0.42-0.08)/(0.42+0.08)=0.34/0.50=0.680$ \ans $f_c=[(0.680-0.15)/(0.85-0.15)]^2=(0.757)^2=0.573$ \ans About 57 per cent of the ground is covered. \finans

}{}
\banner{The theory step by step}
\IfFileExists{extra/th06.tex}{\section{Methods of irrigation and their classification}
The methods of applying water to a field are placed in four groups, by the way in which the water reaches the soil.
\begin{exlist}
\item \kw{Surface irrigation.} Water flows over the soil by gravity and infiltrates as it moves (basin, border, furrow). It needs a levelled field and little energy (Chapter~\ref{chap:11}).
\item \kw{Sprinkler irrigation.} Water is sprayed through the air under pressure and falls like rain (solid-set, hand-move, centre pivot, travelling gun). It suits undulating land and light soils (Chapter~\ref{chap:13}).
\item \kw{Drip (micro) irrigation.} Water is delivered at low pressure through emitters on or near the plant, to a part of the soil volume (Chapter~\ref{chap:12}).
\item \kw{Subsurface irrigation.} The water table is raised or water is delivered below the surface; the root zone is wetted from below.
\end{exlist}
\begin{center}\small\begin{tabular}{@{}p{2.7cm}p{3.2cm}p{3.5cm}p{3.6cm}@{}}\toprule
\hd{Method} & \hd{Typical $E_a$ (per cent)} & \hd{Energy and cost} & \hd{Best suited to}\\\midrule
Surface & 50 to 70 (up to 80 with care) & low pressure; high labour or land-levelling cost & level fields, medium to heavy soils, cereals, rice\\
Sprinkler & 65 to 80 & high pressure; pump cost & undulating land, light soils, wide range of crops\\
Drip & 85 to 95 & low pressure, filters; high capital cost & orchards, vegetables, scarce water, saline water\\\bottomrule
\end{tabular}\end{center}
\noindent The values are typical of well-managed systems; poor management can give much less. Each method is chosen from the water supply, the soil and the slope, the crop, the labour and the capital available.

\section{The efficiencies of an irrigation system}
The water that leaves the source does not all reach the crop, and each stage has a loss. The efficiencies are ratios and are multiplied:
\begin{equation*}\boxed{E_c=\frac{V_{field}}{V_{source}},\qquad E_a=\frac{V_{stored}}{V_{field}},\qquad E_{p}=E_c\,E_a}\end{equation*}
where $E_c$ is the \kw{conveyance efficiency} (the share that reaches the field from the source, reduced by seepage and evaporation in canals), $E_a$ the \kw{application efficiency} (the share of the water delivered to the field that is stored in the root zone), and $E_p$ the \kw{project efficiency}. Two more measures describe the field: the \kw{storage efficiency} (the share of the water needed to refill the root zone that was stored) and the \kw{distribution uniformity} (the evenness of the depth along the field). A system can have a high $E_a$ and a poor uniformity if the deficit is satisfied only in the driest part.

\section{Uniformity}
The depths applied are never perfectly even. The \kw{Christiansen uniformity coefficient} is
\begin{equation*}\boxed{CU=100\left(1-\frac{\sum|x_i-\bar x|}{n\,\bar x}\right)}\end{equation*}
where $x_i$ are the depths collected at $n$ points and $\bar x$ their mean. The \kw{distribution uniformity} of the low quarter is $DU_{lq}=100\,\bar x_{lq}/\bar x$, where $\bar x_{lq}$ is the mean of the lowest quarter of the depths; and for drip systems the \kw{emission uniformity} is based on the lowest quarter of the emitter flows (Chapter~\ref{chap:12}). Good values are $CU>85$ per cent for sprinklers and $DU_{lq}>75$ per cent for surface irrigation.

\section{The leaching requirement}
All irrigation water contains salts, and the crop takes up water and leaves most of the salt in the soil. If the salt is not removed, it builds up in the root zone. The \kw{leaching requirement} is the fraction of the applied water that must pass below the roots to keep the salinity at the level the crop can tolerate:
\begin{equation*}\boxed{LR=\frac{EC_w}{5\,EC_e-EC_w}}\end{equation*}
where $EC_w$ is the electrical conductivity of the irrigation water and $EC_e$ the conductivity of the saturation extract that the crop tolerates (both in dS\,m$^{-1}$). For example, water of $EC_w=1.0$ and a crop with $EC_e=3.0$ give $LR=1.0/(15-1.0)=0.071$: 7 per cent more water than the crop needs must be applied. Drainage must carry this water away (Chapter~\ref{chap:14}).

\section{Weather data: what is needed and why}
The reference evapotranspiration needs five variables: the air temperature (minimum and maximum), the humidity (dew point or relative humidity), the wind speed at a standard height, the radiation (or the hours of sunshine) and the altitude of the station (for the pressure). Rain is needed for the balance. A weather station records them at 2 m height over short grass; data from a distant station should be used with care, because rain varies in space more than the temperature does. When a variable is missing, the FAO gives rules to estimate it (radiation from the range of temperature, humidity from the minimum temperature), at a cost in accuracy.

\section{Vegetation indices from satellite images}
Healthy green leaves reflect strongly in the near infrared and absorb red light, and dry soil reflects both about equally. The \kw{normalised difference vegetation index} is
\begin{equation*}\boxed{NDVI=\frac{\rho_{NIR}-\rho_{red}}{\rho_{NIR}+\rho_{red}}}\end{equation*}
with $\rho$ the reflectance in the band named. It lies between $-1$ and $+1$; bare soil is near 0.1 to 0.2 and a dense green crop is 0.7 to 0.9. The index is related to the green cover and to the crop coefficient, and it is used to map the variation of a field and to divide it into zones. Its weaknesses are saturation at high cover and the effect of the soil background at low cover.
}{}
\banner{Mathematics of irrigation design}
\section{Net and gross requirement}
The net depth $Z_n$ is the water that the root zone must receive to refill the depleted store; the gross depth is $Z_g=Z_n/E_a$ with $E_a$ the application efficiency. The soil and the trigger of the simulation give $Z_n$: for the screenshot of Figure~\ref{fig:sid} the required depth is 72 mm.
\fig{sc_sidk}{0.7}{The surface irrigation layer. The required depth, 72 mm, comes from the simulation (soil, root depth and depletion trigger).}\label{fig:sid}

\section{Surface irrigation: intake and opportunity time}
The cumulative intake of a furrow or border at a point is often described by the Kostiakov relation
\begin{equation}\label{eq:kost}\boxed{Z(\tau)=k\,\tau^{a},}\end{equation}
where $\tau$ is the \emph{opportunity time}, the time during which water is on the soil at that point, and $k,a$ are constants for the soil. The opportunity time needed to infiltrate $Z_n$ is $\tau=(Z_n/k)^{1/a}$.
\exbox{Opportunity time}{With $k=4.0$ mm\,min$^{-a}$ and $a=0.55$ (illustrative values), $Z_n=72$ mm needs $\tau=(72/4.0)^{1/0.55}=18^{1.818}=192$ min. If the inflow per metre width is $q=1.05$ L\,s$^{-1}$\,m$^{-1}$ and the field length is $L=150$ m, then in the ideal case in which every point receives exactly 192 min of water the applied depth is $qt/L=(1.05\times10^{-3})(192\times60)/150=0.0806$ m $=80.6$ mm, so $E_a\le72/80.6=0.89$. Real borders do worse, because water is on the upper end longer than on the lower end (the advance time) and some runs off the tail; reducing this loss is the purpose of the cutoff and cutback rules of the layer.}



\section{Sprinklers: application rate and uniformity}
A sprinkler with discharge $q_s$ on a spacing $S_l\times S_m$ applies depth at the rate
\begin{equation}\label{eq:sprrate}\boxed{r=\frac{q_s}{S_l\,S_m}.}\end{equation}
\par\noindent where $r$ is the application rate (mm\,h$^{-1}$), $q_s$ the discharge of one sprinkler (m$^3$\,h$^{-1}$; multiply the result by 1000 to get mm\,h$^{-1}$), $S_l$ the spacing along the lateral (m) and $S_m$ the spacing between laterals (m).

\exbox{A set time}{With $q_s=1.2$ m$^3$\,h$^{-1}$ on a $12\times18$ m spacing, $r=1.2/216=0.00556$ m\,h$^{-1}=5.56$ mm\,h$^{-1}$. To apply a net depth of 72 mm at $E_a=0.75$ the gross depth is $96$ mm, and the set time is $96/5.56=17.3$ h. The rate must also stay below the intake rate of the soil, or water ponds and runs off: check $r$ against the infiltration rate of Chapter~13 (for the loam, about $\Ks=24.96$ cm\,d$^{-1}=104$ mm\,h$^{-1}$ at long times, far above 5.56 mm\,h$^{-1}$).}


\subhd{Uniformity.} A catch-can test of $N$ collectors gives depths $x_j$ with mean $\bar x$. The Christiansen coefficient is
\begin{equation}\label{eq:cu}\boxed{CU=100\Bigl(1-\frac{\sum_j|x_j-\bar x|}{N\bar x}\Bigr).}\end{equation}
\par\noindent where $CU$ is Christiansen's uniformity coefficient (per cent), $x_j$ the depth caught in can $j$ (mm), $\bar x$ the mean depth and $N$ the number of cans.

\exbox{A catch-can test}{Ten cans collect $8.2,\ 9.1,\ 10.5,\ 7.4,\ 8.8,\ 9.6,\ 11.2,\ 8.0,\ 9.9,\ 10.1$ mm. The mean is $9.28$ mm, the absolute deviations sum to $9.8$ mm, so $CU=100(1-9.8/92.8)=89.4$ per cent. Values above about 85 per cent are usually considered acceptable for field crops.}


\fig{sc_sprk}{0.7}{The SWAFL sprinkler layer: nozzle, spacing, wind and soil intake.}

\section{Drip: rate per area}
A drip line has emitters of discharge $q_e$ at spacing $s_e$ on laterals at spacing $s_l$. The depth rate over the whole area is $r=q_e/(s_e s_l)$.
\exbox{A drip schedule}{With $q_e=2$ L\,h$^{-1}$, $s_e=0.5$ m and $s_l=1.5$ m, $r=2/(0.5\times1.5)=2.67$ mm\,h$^{-1}$. A daily gross requirement of 6 mm needs $6/2.67=2.25$ h. The wetted bulb below an emitter is narrow and deep in sand and wide and shallow in clay, because the capillary pull (Chapter~3) is weak in a sand, where $K$ at the wet end is large and gravity dominates.}


\fig{sc_dripk}{0.7}{The SWAFL drip layer: emitters, wetted bulb and uniformity.}

\subhd{Pipe network.}
The friction loss of Equation~(\ref{eq:hw}) and the speed $v=4Q/(\pi D^2)$ are computed for every pipe of the sketch, and the pump head is the static lift plus the largest sum of losses along a path (including the operating pressure of the emitters or sprinklers).

\section{Subsurface drainage: the Hooghoudt equation}
For parallel drains at spacing $L$ above an impermeable layer, the steady-state water table midway between the drains has a height $h$ above drain level given by the Hooghoudt equation,
\begin{equation}\label{eq:hooghoudt}\boxed{L^{2}=\frac{8K_2\,d_e\,h+4K_1h^{2}}{q},}\end{equation}
where $q$ is the drainage rate (m\,d$^{-1}$), $K_1$ and $K_2$ the conductivities above and below drain level, and $d_e$ the equivalent depth that accounts for convergence of the flow towards the drain.
\exbox{Drain spacing}{With $K_1=K_2=0.5$ m\,d$^{-1}$, $d_e=1.5$ m, $h=0.60$ m and $q=0.005$ m\,d$^{-1}$: $L^2=(8\cdot0.5\cdot1.5\cdot0.6+4\cdot0.5\cdot0.36)/0.005=(3.6+0.72)/0.005=864$ m$^2$, so $L=29.4$ m. Halving the conductivity to 0.25 m\,d$^{-1}$ halves $L^2$ to 432, so $L=20.8$ m.}


\fig{sc_drnk}{0.7}{The SWAFL drainage layer.}

\exbox{}{With $k=3.0$ and $a=0.6$ find the opportunity time for $Z_n=60$ mm.}\sol \givenline{$k=3.0$ ; $a=0.6$ find ; $Z_n=60$ mm}\ans \[\tau=(60/3)^{1/0.6}=20^{1.667}=147\ \text{min}\quad\text{\finans}\]

\exbox{}{A sprinkler of 2.0 m$^3$\,h$^{-1}$ is on a $15\times15$ m spacing. Find $r$ and the set time for a gross depth of 80 mm.}\sol  $r=2.0/225=8.89$ mm\,h$^{-1}$; $80/8.89=9.0$ h. \finans

\exbox{}{Compute $CU$ for cans of $6,\ 8,\ 9,\ 10,\ 12$ mm.}\sol  Mean $9$; deviations $3,1,0,1,3$ sum to $8$; $CU=100(1-8/45)=82.2$ per cent. \finans

\exbox{}{A drip system has $q_e=4$ L\,h$^{-1}$ at 0.75 m spacing on laterals 2 m apart. Find $r$ and the time to apply 5 mm gross.}\sol \givenline{$q_e=4$ L\,h$^{-1}$}\ans $r=4/(0.75\cdot2)=2.67$ mm\,h$^{-1}$; $5/2.67=1.88$ h. \finans

\exbox{}{Find the Hooghoudt spacing for $K=1$ m\,d$^{-1}$ ($K_1=K_2$), $d_e=2$ m, $h=0.5$ m and $q=0.007$ m\,d$^{-1}$.}\sol \givenline{$K=1$ m\,d$^{-1}$ ; $K_1=K_2$ ; $d_e=2$ m ; $h=0.5$ m ; $q=0.007$ m\,d$^{-1}$}\ans \[L^2=(8\cdot1\cdot2\cdot0.5+4\cdot1\cdot0.25)/0.007=9/0.007=1{,}286\]
$L=35.9$ m. \finans

\exbox{}{Why can $E_a$ of Equation~(\ref{eq:kost})'s ideal case not be reached in a real border?}\sol  Water arrives at the top before the bottom, so the top gets a longer opportunity time than the bottom, and some water runs off the tail.
\banner{Weather and field data}\label{ch:weather}
\section{Reference evapotranspiration}
Without an ET$_0$ column in the weather file the program computes it. With daily maximum and minimum temperature it uses Hargreaves--Samani (FAO-56 Eq.~52),
\begin{equation}\boxed{ET_0=0.0023(T_{\rm mean}+17.8)\sqrt{T_{\max}-T_{\min}}\;R_a\times0.408\ \ (\mathrm{mm\,d^{-1}}),}\end{equation}
where $R_a$ (MJ m$^{-2}$ d$^{-1}$) is the extraterrestrial radiation at the latitude typed in the dialog (FAO-56 Eq.~21),
\begin{equation}R_a=\frac{24\cdot60}{\pi}G_{sc}d_r\left[\omega_s\sin\varphi\sin\delta+\cos\varphi\cos\delta\sin\omega_s\right],\ d_r=1+0.033\cos\frac{2\pi J}{365},\ \delta=0.409\sin\left(\frac{2\pi J}{365}-1.39\right),\ \omega_s=\arccos(-\tan\varphi\tan\delta),\end{equation}
\par\noindent where $R_a$ is the extraterrestrial radiation (MJ\,m$^{-2}$\,d$^{-1}$), $G_{sc}$ the solar constant, $d_r$ the relative distance of the earth from the sun, $\delta$ the solar declination (rad), $\omega_s$ the sunset hour angle (rad), $\varphi$ the latitude (rad) and $J$ the day of the year. \finans

$G_{sc}=0.0820$ MJ m$^{-2}$ min$^{-1}$. With the daily mean temperature and the global radiation $R_s$ (MJ m$^{-2}$ d$^{-1}$) it uses Makkink,
\begin{equation}\boxed{ET_0=0.61\frac{\Delta}{\Delta+\gamma}\frac{R_s}{2.45}-0.12,\quad\Delta=\frac{4098\times0.6108\,e^{17.27T/(T+237.3)}}{(T+237.3)^2},\ \gamma=0.066\ \mathrm{kPa\,K^{-1}}.}\end{equation}
\par\noindent where $ET_0$ is the reference evapotranspiration (mm\,d$^{-1}$) estimated from the solar radiation $R_s$ (MJ\,m$^{-2}$\,d$^{-1}$), 2.45 is the latent heat of vaporisation (MJ\,kg$^{-1}$), $\Delta$ the slope of the vapour-pressure curve (kPa\,K$^{-1}$), $T$ the air temperature (\textdegree C) and $\gamma$ the psychrometric constant.


\exbox{Hargreaves at 10$^\circ$N}{A site at latitude $10^\circ$N on day of year $J=180$ has $T_{\max}=33^\circ$C and $T_{\min}=21^\circ$C. Compute $R_a$ and Hargreaves ET$_0$.}\sol \givenline{$J=180$ has ; $T_{\max}=33^\circ$ C ; $T_{\min}=21^\circ$ C}\ans $\varphi=10^\circ=0.1745$ rad. \[d_r=1+0.033\cos(2\pi\times180/365)=1+0.033\times-0.9991=0.9670\]
\[\delta=0.409\sin(2\pi\times180/365-1.39)=0.409\sin(1.7086)=0.4051\] rad. \[\omega_s=\arccos(-\tan\varphi\tan\delta)=\arccos(-0.1763\times0.4288)=1.6465\] rad. Then \[R_a=(1440/\pi)\times0.0820\times0.9670\times[1.6465\times\sin\varphi\sin\delta+\cos\varphi\cos\delta\sin\omega_s]=36.90\ \text{MJ m$^{-2}$ d$^{-1}$}\]
$T_{\rm mean}=27^\circ$C, $\sqrt{33-21}=3.464$. \[ET_0=0.0023\times(27+17.8)\times3.464\times36.90\times0.408=5.37\ \text{mm d$^{-1}$}\quad\text{\finans}\]


\begin{table}[!htbp]\centering\small\caption{Extraterrestrial radiation (MJ m$^{-2}$ d$^{-1}$) by latitude and day, and Hargreaves ET$_0$ at midsummer of the northern hemisphere (mm d$^{-1}$).}\label{tab:ra}
\begin{tabular}{lrrrrr}\toprule
\hd{Latitude} & \hd{$R_a$ J=15} & \hd{J=105} & \hd{J=196} & \hd{J=288} & \hd{Hargreaves J=196 (30/18$^\circ$C)}\\\midrule
0 & 36.2 & 36.8 & 33.9 & 37.3 & 4.60\\
10 & 31.9 & 37.9 & 37.1 & 35.1 & 5.03\\
30 & 21.1 & 36.8 & 40.5 & 27.5 & 5.51\\
50 & 8.9 & 31.5 & 40.2 & 16.8 & 5.46\\\bottomrule\end{tabular}\end{table}

\exbox{Makkink}{Makkink: $T=27^\circ$C, $R_s=20$ MJ m$^{-2}$ d$^{-1}$.}\sol \givenline{$T=27^\circ$ C ; $R_s=20$ MJ m$^{-2}$ d$^{-1}$}\ans \[e^{17.27\times27/(27+237.3)}=e^{1.7642}=5.8372\]
saturation pressure $0.6108\times5.8372=3.5653$ kPa; \[\Delta=4098\times3.5653/(264.3)^2=0.2092\ \text{kPa K$^{-1}$}\]
\[\Delta/(\Delta+\gamma)=0.2092/(0.2092+0.066)=0.7601\]
\[ET_0=0.61\times0.7601\times20/2.45-0.12=3.67\ \text{mm d$^{-1}$}\quad\text{\finans}\]


\subsection*{The FAO Penman--Monteith value of the online weather}
When the weather of a place is fetched the reference ET$_0$ comes from the weather service (FAO Penman--Monteith). The equation is the literature one (FAO-56 Eq.~6):
\begin{equation}\boxed{ET_0=\frac{0.408\Delta(R_n-G)+\gamma\frac{900}{T+273}u_2(e_s-e_a)}{\Delta+\gamma(1+0.34u_2)},\quad G\approx0\ \text{(daily)}.}\end{equation}
\par\noindent where Symbols are as in Eq.~(\ref{eq:pm}); $G$ is the soil heat flux, taken as zero for a daily step.

The program does not compute this itself (the service does), so the example below is for understanding and for checking the Hargreaves value, not a reproduction of the app.

\exbox{FAO Penman--Monteith for comparison}{Same day: $T_{\max}=33$, $T_{\min}=21^\circ$C, mean RH 60\%, wind $u_2=2.0$ m s$^{-1}$, $R_s=20$ MJ m$^{-2}$ d$^{-1}$, altitude 100 m.}\sol \givenline{$T_{\max}=33$ ; $T_{\min}=21^\circ$ C ; $u_2=2.0$ m s$^{-1}$ ; $R_s=20$ MJ m$^{-2}$ d$^{-1}$}\ans $T=27$; pressure \[P=101.3[(293-0.0065\times100)/293]^{5.26}=100.12\ \text{kPa}\]
$\gamma=0.000665P=0.06658$. \[e_s=[e^\circ(33)+e^\circ(21)]/2=3.759\ \text{kPa}\]
$e_a=0.60e_s=2.255$. $\Delta=0.2092$. $R_a=36.90$; clear-sky \[R_{so}=(0.75+2\times10^{-5}\times100)R_a=27.75\]
$R_{ns}=0.77\times20=15.40$; \[R_{nl}=4.903\times10^{-9}[\ldots](0.34-0.14\sqrt{e_a})(1.35R_s/R_{so}-0.35)=3.23\]
$R_n=12.17$. Then \[ET_0=[0.408\times0.2092\times12.17+0.06658\times900/300\times2\times(3.759-2.255)]/[0.2092+0.06658(1+0.68)]=5.11\ \text{mm d$^{-1}$}\]
Compare Hargreaves (5.37) and Makkink (3.67): the three differ by up to about 1.7 mm d$^{-1}$, which is why the layer reports how ET$_0$ was obtained. \finans


\section{Reading and gridding a weather file}
Radiation columns are read in MJ m$^{-2}$ d$^{-1}$, or in W m$^{-2}$ when the values exceed 60 (a daily mean power: multiply by $0.0864$); in the Daisy files radiation is in kWh m$^{-2}$ d$^{-1}$ ($\times3.6$ gives MJ m$^{-2}$ d$^{-1}$). Rain in inches is multiplied by 25.4. The program works on twenty years of 365 days: the last twenty complete years (at least 300 days each) are used, leap days dropped, a missing day of rain taken from the same day of another year, a missing ET$_0$ from the mean of that day of the year.

\exbox{Units}{A file gives mean radiation 230 W m$^{-2}$, another 5.6 kWh m$^{-2}$ d$^{-1}$, a third 18.4 MJ m$^{-2}$ d$^{-1}$, and a rain value of 0.8 in. Convert all to model units.}\sol  $230>60$, so it is W m$^{-2}$: $230\times0.0864=19.87$ MJ m$^{-2}$ d$^{-1}$. $5.6\times3.6=20.16$ MJ m$^{-2}$ d$^{-1}$. $18.4$ is already in MJ. Rain: $0.8\times25.4=20.3$ mm. Note the threshold: a value of 55 would be read as MJ, which is unphysical (the extraterrestrial maximum is about 45), so the 60 limit is safe for daily totals; a mean power of 55 W m$^{-2}$ would be misread, a limitation of the rule. \finans


The daily model then uses the rain of the file scaled to the rainfall scale $\rho$ of the Properties panel and a demand weight proportional to the daily ET$_0$:
\begin{equation}P_k^{\rm model}=P_k\frac{\rho}{\rho_{\rm ref}},\qquad e_k=\frac{ET_{0,k}}{\overline{ET_0}/365},\qquad\rho_{\rm ref}=\mathop{\mathrm{nearest}}_{\rho\in\{0.5,\dots,1.4\}}\frac{\bar P}{P_{\rm rec}},\qquad\Delta T_{\rm ref}=\mathop{\mathrm{nearest}}_{\Delta T\in\{-1,\dots,3.5\}}\frac{\overline{ET_0}/E_{\rm rec}-1}{0.035},\end{equation}
with $P_{\rm rec}=464$ mm the mean annual rain of the built-in record, and $E_{\rm rec}\approx870$ mm its reference ET. The annual model (Engine~1) uses the simulated steps nearest the weather.

\exbox{Placing a file on the grid}{A weather file has mean annual rain $\bar P=640$ mm and $\overline{ET_0}=1{,}150$ mm. Find $\rho_{\rm ref}$ and $\Delta T_{\rm ref}$, and say whether the annual model can represent the file.}\sol \givenline{$\bar P=640$ mm ; $\overline{ET_0}=1{,}150$ mm}\ans \[\bar P/P_{\rm rec}=640/464=1.379\]
The nearest step of $\{0.5,0.6,\dots,1.4\}$ is $1.4$. \[\overline{ET_0}/E_{\rm rec}=1150/870=1.322\]
$(1.322-1)/0.035=9.2$ $^\circ$C. The warming steps stop at 3.5, so the nearest is $3.5$ and the file lies outside the simulated range: the annual results use the nearest step and the dialog says so. The daily layers use the file as it is. A second file with $\bar P=900$ mm gives $900/464=1.940$ and is also capped at 1.4. \finans


\section{Field geometry}
The field outline $(\varphi_i,\lambda_i)$ (radians) gives the area and the perimeter on a sphere of radius $R=6\,378\,137$ m,
\begin{equation}A=\frac{R^2}2\Big|\sum_i(\lambda_{i+1}-\lambda_i)(2+\sin\varphi_i+\sin\varphi_{i+1})\Big|,\qquad P=\sum_i2R\arcsin\sqrt{\sin^2\tfrac{\Delta\varphi}{2}+\cos\varphi_i\cos\varphi_{i+1}\sin^2\tfrac{\Delta\lambda}{2}}.\end{equation}
\par\noindent where $A$ is the area of the polygon (m$^2$), $P$ its perimeter (m), $R$ the radius of the earth (6\,371\,000 m), $\lambda_i$ and $\varphi_i$ the longitude and latitude of corner $i$ (rad), $\Delta\varphi=\varphi_{i+1}-\varphi_i$ and $\Delta\lambda=\lambda_{i+1}-\lambda_i$.


\exbox{Area of a field from its corners}{A rectangular field has corners at $9.000^\circ$N, $7.500^\circ$E and extends about 300 m north and 400 m east. Compute the area and the perimeter from the two formulas above.}\sol  The north--south side is $300/111{,}320=0.002695^\circ$; the east--west side \[400/(111{,}320\cos9^\circ)=0.003638^\circ\]
Summing the four terms of the area formula gives $A=119{,}998$ m$^2=12.000$ ha. The haversine perimeter is $1{,}399.99$ m, within 0.01 m of the plane value $2(300+400)=1{,}400$ m, and $A$ is within 2 m$^2$ of $300\times400=120{,}000$ m$^2$. At this size the curvature of the Earth does not matter; the spherical formulas are there so that large or distant outlines are right, and this agreement is the check. \finans


\subsection*{Slope by a fitted plane}
The ground slope is that of the least-squares plane $z=a+bx+cy$ through the elevations ($x$, $y$ in metres east and north of the centre): $s=100\sqrt{b^2+c^2}$ per cent, falling in the direction of $-(b,c)$.

\begin{table}[!htbp]\centering\small\caption{Nine elevations over the field.}\label{tab:dem}
\begin{tabular}{lrr}\toprule
\hd{$x$ (m)} & \hd{$y$ (m)} & \hd{$z$ (m)}\\\midrule
-150 & -100 & 512.3\\
0 & -100 & 511.2\\
150 & -100 & 510.4\\
-150 & 0 & 512.9\\
0 & 0 & 511.6\\
150 & 0 & 510.7\\
-150 & 100 & 513.4\\
0 & 100 & 512.2\\
150 & 100 & 511.0\\\bottomrule\end{tabular}\end{table}

\exbox{Slope of the field}{Fit the plane to Table~\ref{tab:dem}; give the slope in per cent and the downhill direction.}\sol  The nine points are symmetric about the centre, so the normal equations decouple: $a=\bar z=511.744$, \[b=\sum xz/\sum x^2=-0.00722\]
\[c=\sum yz/\sum y^2=0.00450\] (\[\sum x^2=6\times150^2=135{,}000\]
\[\sum y^2=6\times100^2=60{,}000\]). \[s=100\sqrt{-0.00722^2+0.00450^2}=0.85\%\]
The ground falls toward $-(b,c)$, that is toward the east-south-east: bearing $122^\circ$ from north. A slope of 0.85\% is the value used by the sprinkler layer (Chapter~\ref{ch:sprinkler}) to look up the allowed application rate (0--5\% column). \finans


\subsection*{Where to take the water from}
\begin{equation}\boxed{\mathrm{score}=B_{\rm type}-30[\text{seasonal}]-\min(45,d/60)-\min(20,0.8\Delta z),}\end{equation}
with $B=92$ river, 90 canal, 86 reservoir, 82 lake, 74 borehole, 70 well, 62 pond, 60 spring, 58 stream, 50 water seen in the photo, 42 ditch, 30 drain; $d$ the distance from the field edge (m), $\Delta z$ the lift (m). The weights are engineered, not fitted.

\begin{table}[!htbp]\centering\small\caption{Five candidate sources.}\label{tab:src}
\begin{tabular}{lrrrrr}\toprule
\hd{Source} & \hd{$B$} & \hd{seasonal} & \hd{$d$ (m)} & \hd{lift (m)} & \hd{score}\\\midrule
River (seasonal) & 92 & yes & 900 & 8 & 40.6\\
Canal & 90 & no & 1200 & 3 & 67.6\\
Reservoir & 86 & no & 400 & 25 & 59.3\\
Borehole & 74 & no & 150 & 0 & 71.5\\
Pond & 62 & no & 100 & 4 & 57.1\\\bottomrule\end{tabular}\end{table}

\exbox{Choosing the water source}{Score the sources of Table~\ref{tab:src}.}\sol  River: \[92-30-\min(45,900/60)-\min(20,0.8\times8)=92-30-15-6.4=40.6\]
Canal: $90-0-20-2.4=67.6$. Reservoir: $86-6.67-20=59.3$ (the lift term is capped at 20). Borehole: $74-2.5-0=71.5$. Pond: $62-1.67-3.2=57.1$. The best is the Borehole ($71.5$). The seasonal river drops below the canal because of its $-30$ penalty even though its type score is higher. \finans


\section{Vegetation and water indices}
For each 10 m cell with blue, red, red-edge, NIR and SWIR reflectances $\rho_B,\rho_R,\rho_{RE},\rho_N,\rho_{SW}$:
\begin{equation}\mathrm{NDVI}=\frac{\rho_N-\rho_R}{\rho_N+\rho_R},\ \mathrm{EVI}=2.5\frac{\rho_N-\rho_R}{\rho_N+6\rho_R-7.5\rho_B+1},\ \mathrm{SAVI}=\frac{(1+L)(\rho_N-\rho_R)}{\rho_N+\rho_R+L},\ L=0.5,\ \mathrm{NDRE}=\frac{\rho_N-\rho_{RE}}{\rho_N+\rho_{RE}},\ \mathrm{NDWI}=\frac{\rho_N-\rho_{SW}}{\rho_N+\rho_{SW}}.\end{equation}
\par\noindent where $\rho_N$, $\rho_R$, $\rho_B$, $\rho_{RE}$ and $\rho_{SW}$ are the reflectances in the near-infrared, red, blue, red-edge and short-wave infrared bands (0 to 1), $L$ the soil-adjustment factor and NDVI, EVI, SAVI, NDRE, NDWI the resulting vegetation and water indices.

From the colour photo only: $\mathrm{VARI}=\frac{G-R}{G+R-B}$, $\mathrm{GLI}=\frac{2G-R-B}{2G+R+B}$, $\mathrm{ExG}=2g-r-b$ (chromatic coordinates), $\mathrm{NGRDI}=\frac{G-R}{G+R}$.

\begin{table}[!htbp]\centering\small\caption{Indices from illustrative reflectances (blue, red, red-edge, NIR, SWIR): dense (0.03, 0.04, 0.18, 0.45, 0.20), sparse (0.05, 0.10, 0.20, 0.30, 0.25), bare (0.09, 0.20, 0.23, 0.27, 0.35), stressed (0.04, 0.08, 0.20, 0.32, 0.28).}\label{tab:idx}
\begin{tabular}{lrrrrr}\toprule
\hd{Cell type} & \hd{NDVI} & \hd{EVI} & \hd{SAVI} & \hd{NDRE} & \hd{NDWI}\\\midrule
Dense crop & 0.837 & 0.700 & 0.621 & 0.429 & 0.385\\
Sparse crop & 0.500 & 0.328 & 0.333 & 0.200 & 0.091\\
Bare soil & 0.149 & 0.097 & 0.108 & 0.080 & -0.129\\
Stressed crop & 0.600 & 0.400 & 0.400 & 0.231 & 0.067\\\bottomrule\end{tabular}\end{table}

\exbox{Indices of one cell}{Compute the indices of the dense-crop cell $(\rho_B,\rho_R,\rho_{RE},\rho_N,\rho_{SW})=(0.03,0.04,0.18,0.45,0.20)$.}\sol  NDVI \[=(0.45-0.04)/(0.45+0.04)=0.41/0.49=0.837\]
EVI \[=2.5\times0.41/(0.45+0.24-0.225+1)=1.025/1.465=0.700\]
SAVI \[=1.5\times0.41/(0.45+0.04+0.5)=0.615/0.99=0.621\]
NDRE \[=(0.45-0.18)/(0.45+0.18)=0.27/0.63=0.429\]
NDWI \[=(0.45-0.20)/(0.45+0.20)=0.25/0.65=0.385\]
For the bare cell NDVI is only 0.149, and its NDWI is negative (-0.129): bare soil reflects more in the SWIR than in the NIR. \finans


\exbox{Visible-band indices}{A drone photo cell has $R=110$, $G=150$, $B=60$ (8-bit). Compute VARI, GLI, ExG and NGRDI.}\sol \givenline{$R=110$ ; $G=150$ ; $B=60$}\ans Sum $R+G+B=320$; $r=0.3438$, $g=0.4688$, $b=0.1875$. VARI \[=(150-110)/(150+110-60)=40/200=0.200\]
GLI \[=(300-110-60)/(300+110+60)=130/470=0.277\]
ExG \[=2\times0.4688-0.3438-0.1875=0.406\]
NGRDI $=40/260=0.154$. All four are positive: the cell is greener than red. \finans


\section{Zones, cover and the zone-by-zone balance}
From the series of scenes each cell gets the standardised index $z_{i,t}=(\mathrm{NDVI}_{i,t}-\overline{\mathrm{NDVI}}_t)/\sigma_t$, the stability $\frac1n\sum_tz_{i,t}$ and the change $\Delta_i=\mathrm{NDVI}_{i,t_n}-\mathrm{NDVI}_{i,t_{n-1}}$. For the zone balance the cover is $f_c=(\mathrm{NDVI}-0.15)/0.7$, scaled by the model's crop development $c(t)$ as $f_{c,z}(t)=\min(1,c(t)f_c(\mathrm{NDVI}_z)/\bar f_c)$, and the crop coefficient $K_c=K_{c,\rm ini}+(K_{c,\rm mid}-K_{c,\rm ini})K_d$ with $K_d=\min(1,1.5f_c,f_c^{0.5})$ (Allen and Pereira, 2009).

\begin{table}[!htbp]\centering\small\caption{NDVI of four cells over five scenes; population standard deviation.}\label{tab:ndvi}
\begin{tabular}{lrrrrrrr}\toprule
 & \hd{t1} & \hd{t2} & \hd{t3} & \hd{t4} & \hd{t5} & \hd{stability} & \hd{$\Delta$}\\\midrule
cell 1 & 0.42 & 0.55 & 0.66 & 0.71 & 0.62 & 0.30 & -0.09\\
cell 2 & 0.40 & 0.57 & 0.70 & 0.74 & 0.60 & 0.41 & -0.14\\
cell 3 & 0.30 & 0.44 & 0.52 & 0.55 & 0.47 & -1.66 & -0.08\\
cell 4 & 0.45 & 0.60 & 0.72 & 0.76 & 0.66 & 0.95 & -0.10\\
field mean & 0.393 & 0.540 & 0.650 & 0.690 & 0.588 &  & \\
$\sigma_t$ & 0.056 & 0.060 & 0.078 & 0.083 & 0.071 &  & \\\bottomrule\end{tabular}\end{table}

\exbox{Stability of a cell}{For Table~\ref{tab:ndvi} find $z_{3,t}$, the stability of cell 3 and its change since the previous scene, and say which cell belongs in the weakest monitoring zone.}\sol  Field means $(0.393, 0.540, 0.650, 0.690, 0.588)$; standard deviations $(0.056, 0.060, 0.078, 0.083, 0.071)$. For cell 3: \[z=(0.30-0.393)/0.056=-1.64\]
then -1.66, -1.66, -1.69, -1.65. Stability $=\frac15\sum z=-1.66$; $\Delta_3=0.47-0.55=-0.08$. Cell 3 is below the field on every date (stability $-1.66$), so it belongs to the weakest zone: a zone is a part of the field that is weaker on every date, not on one. \finans


\begin{table}[!htbp]\centering\small\caption{Cover and crop coefficient from NDVI ($K_{c,\rm ini}=0.30$, $K_{c,\rm mid}=1.20$).}\label{tab:ndvikc}
\begin{tabular}{lrrr}\toprule
\hd{NDVI} & \hd{$f_c$} & \hd{$K_d$} & \hd{$K_c$ (maize)}\\\midrule
0.20 & 0.071 & 0.107 & 0.396\\
0.30 & 0.214 & 0.321 & 0.589\\
0.45 & 0.429 & 0.643 & 0.879\\
0.60 & 0.643 & 0.802 & 1.022\\
0.75 & 0.857 & 0.926 & 1.133\\
0.85 & 1.000 & 1.000 & 1.200\\\bottomrule\end{tabular}\end{table}

\exbox{Crop coefficient from NDVI}{A zone has NDVI = 0.60. Find $f_c$, $K_d$ and the maize $K_c$.}\sol  $f_c=(0.60-0.15)/0.7=0.643$. \[K_d=\min(1,1.5\times0.643,\sqrt{0.643})=\min(1,0.964,0.802)=0.802\]
\[K_c=0.30+0.90\times0.802=1.022\]
For a thin canopy (NDVI 0.30), $f_c=0.214$, \[K_d=\min(1,0.321,0.463)=0.321\] and $K_c=0.589$: $K_d$ rises faster than the cover at first (the $1.5f_c$ branch) and then follows $\sqrt{f_c}$. \finans


\section{Runoff, gullies and soil loss}
\textbf{Runoff.} The runoff of each day follows the SCS curve number (TR-55): with $S=25400/CN-254$ (mm) and $I_a=0.2S$, $Q=(P-I_a)^2/(P+0.8S)$ for $P>I_a$. The curve number depends on the hydrologic soil group (from the SoilGrids texture) and the cover: row crop in good condition, or bare soil when the monthly NDVI is below 0.3. The values below are the standard TR-55 ones for those two classes and are quoted for the worked example; they were not checked against the program.

\begin{table}[!htbp]\centering\small\caption{Runoff $Q$ (mm) from daily rain $P$ = 20, 40, 80 mm.}\label{tab:cn}
\begin{tabular}{lrrrrrrrr}\toprule
\hd{Group} & \hd{CN row} & \hd{CN bare} & \hd{Q 20} & \hd{Q 40} & \hd{Q 80} & \hd{bare Q 20} & \hd{bare Q 40} & \hd{bare Q 80}\\\midrule
A & 67 & 77 & 0.0 & 1.6 & 16.8 & 0.3 & 6.1 & 29.9\\
B & 78 & 86 & 0.4 & 6.8 & 31.4 & 2.6 & 13.8 & 45.5\\
C & 85 & 91 & 2.2 & 12.7 & 43.6 & 5.6 & 20.4 & 56.2\\
D & 89 & 94 & 4.2 & 17.5 & 51.7 & 8.5 & 25.5 & 63.4\\\bottomrule\end{tabular}\end{table}

\exbox{Runoff of one storm}{A 40 mm rain falls on a group B soil under a row crop ($CN=78$) and then on the same soil bare ($CN=86$).}\sol \givenline{$CN=78$ ; $CN=86$}\ans Row crop: $S=25400/78-254=71.6$ mm, $I_a=0.2S=14.3$ mm, \[Q=(40-14.3)^2/(40+0.8\times71.6)=6.8\ \text{mm}\]
Bare ($CN=86$): $S=41.3$, $I_a=8.3$, $Q=13.8$ mm. Bare soil sheds 2.03 times as much runoff. \finans


\subhd{Ephemeral gullies.} A gully is likely where $G=S\,A^{0.4}/0.025\ge1$, with $A$ the upslope area (ha) and $S$ the slope (m m$^{-1}$) (Vandaele et al., 1996).
\textbf{Soil loss (RUSLE).} $A_s=RKLS\,CP$ with $R=0.0483P_a^{1.61}$ (for $P_a\le850$ mm; $587.8-1.219P_a+0.004105P_a^2$ above), shared over the months by the erosive daily rains ($P_d\ge12.7$ mm, exponent 1.81), $C_i=\sum_m(R_m/R)\,e^{-2\mathrm{NDVI}_{i,m}/(1-\mathrm{NDVI}_{i,m})}$ and $LS=(A/22.13)^{0.4}(\sin\beta/0.0896)^{1.3}$ (Moore and Wilson, 1992; $A$ the specific upslope area).

\begin{table}[!htbp]\centering\small\caption{Gully threshold.}\label{tab:gully}
\begin{tabular}{lrrr}\toprule
\hd{$A$ (ha)} & \hd{$S$ (\%)} & \hd{$G$} & \hd{gully}\\\midrule
0.1 & 1 & 0.16 & no\\
0.1 & 3 & 0.48 & no\\
0.5 & 1 & 0.30 & no\\
0.5 & 3 & 0.91 & no\\
1.0 & 1 & 0.40 & no\\
1.0 & 3 & 1.20 & likely\\
2.0 & 1 & 0.53 & no\\
2.0 & 3 & 1.58 & likely\\
5.0 & 1 & 0.76 & no\\
5.0 & 3 & 2.28 & likely\\\bottomrule\end{tabular}\end{table}

\begin{table}[!htbp]\centering\small\caption{Months of an illustrative year: erosive-rain share, NDVI and the cover factor $C_m=e^{-2\mathrm{NDVI}/(1-\mathrm{NDVI})}$.}\label{tab:rusle}
\begin{tabular}{lrrrr}\toprule
\hd{month} & \hd{erosive share (\%)} & \hd{NDVI} & \hd{$C_m$} & \hd{$R_m$}\\\midrule
1 & 0 & 0.25 & 0.513 & 0\\
2 & 0 & 0.25 & 0.513 & 0\\
3 & 1 & 0.25 & 0.513 & 16\\
4 & 6 & 0.30 & 0.424 & 96\\
5 & 14 & 0.45 & 0.195 & 223\\
6 & 20 & 0.65 & 0.024 & 318\\
7 & 23 & 0.72 & 0.006 & 366\\
8 & 19 & 0.70 & 0.009 & 302\\
9 & 13 & 0.55 & 0.087 & 207\\
10 & 4 & 0.35 & 0.341 & 64\\
11 & 0 & 0.25 & 0.513 & 0\\
12 & 0 & 0.25 & 0.513 & 0\\\bottomrule\end{tabular}\end{table}

\exbox{Soil loss of a field}{Annual rain 640 mm; loam (40\% sand, 40\% silt, 20\% clay, 1.2\% organic carbon); slope 4\% with a specific upslope area of 60 m; monthly NDVI and erosive-rain share of Table~\ref{tab:rusle}. Find the soil loss with $P=1$.}\sol \givenline{$P=1$}\ans \[R=0.0483\times640^{1.61}=1{,}592\] MJ mm ha$^{-1}$ h$^{-1}$ y$^{-1}$. $K$ (EPIC): \[0.2+0.3e^{-0.0256\times40(1-0.40)}=0.3623\]
$[40/60]^{0.3}=0.8855$; organic-carbon term $0.8749$; sand term $0.9999$; $K=0.2806$ (US units), $0.0370$ in t ha h ha$^{-1}$ MJ$^{-1}$ mm$^{-1}$. \[LS=(60/22.13)^{0.4}(\sin\beta/0.0896)^{1.3}\] with $\beta=\arctan0.04$: $\sin\beta=0.0400$, $LS=0.522$. $C=\sum(R_m/R)C_m=0.091$ (the erosive months have a green canopy, so the effective $C$ is much smaller than the bare-soil value $e^{-0.667}=0.51$). So \[A_s=1{,}592\times0.0370\times0.522\times0.091\times1=2.79\] t ha$^{-1}$ y$^{-1}$, above the tolerable 2 t ha$^{-1}$ y$^{-1}$. \finans


\begin{keybox}The soil-loss figure is as good as its factors: the $R$ law was fitted to rain regimes quite different from many sites, the $K$ formula is the EPIC approximation of the nomograph, and the $C$ factor from NDVI is an empirical curve. The layer presents the result against the tolerable 2 t ha$^{-1}$ y$^{-1}$ and names the risk window (months when NDVI is below 0.35 meets a large erosive share); it does not predict a gully's depth.\end{keybox}

\IfFileExists{extra/plus06.tex}{\banner{Further analysis: the energy line of a pipe and the intake curve}
\section{The energy grade line of a pipe}\label{sec:dmegl}
The total head of water in a pipe is the sum of the elevation $z$, the pressure head $p/(\rho_wg)$ and the velocity head $v^2/(2g)$. Along the pipe it falls by the friction loss $h_f$ and by local losses at fittings. Drawing the \kw{energy grade line} and the \kw{hydraulic grade line} (energy line less the velocity head) shows at once where the pressure is lowest, which is where an emitter or a sprinkler will deliver least.
\begin{equation}\label{eq:dmbern}\boxed{z_1+\frac{p_1}{\rho_wg}+\frac{v_1^2}{2g}=z_2+\frac{p_2}{\rho_wg}+\frac{v_2^2}{2g}+h_f}\end{equation}

\begin{kfig}{Energy grade line (EGL) and hydraulic grade line (HGL) of a pipe: the energy falls along the pipe by the friction loss $h_f$; the pressure head is the height of the HGL above the pipe.}
\draw[ln,line width=1.2pt] (0,0.3)--(7,-0.4);\node[lab,below] at (3.5,-0.15) {pipe};
\draw[crv] (0,3.4)--(7,2.0);\draw[gd] (0,3.0)--(7,1.6);
\node[lab,right] at (7,2.0) {EGL};\node[lab,right] at (7,1.6) {HGL};
\draw[<->,ln] (1.5,0.15)--(1.5,2.7);\node[lab,right] at (1.5,1.5) {$p/\rho_wg$};
\draw[<->,ln] (2.5,2.5)--(2.5,2.9);\node[lab,right] at (2.55,2.75) {$v^2/2g$};
\draw[ln] (0,3.4)--(7,3.4);\draw[<->,ln] (6.5,3.4)--(6.5,2.1);\node[lab,left] at (6.45,2.8) {$h_f$};
\draw[ln] (-0.3,-0.6)--(7.3,-0.6);\node[lab,below] at (0.3,-0.6) {datum};
\end{kfig}

\section{The intake curve of Kostiakov}\label{sec:dmkost}
The cumulative infiltration of a soil after an opportunity time $\tau$ is often fitted by $Z=k\tau^{a}$ ($0<a<1$). It answers the question that every surface design asks: how long must water stand on the soil to infiltrate the required depth?
\begin{equation}\label{eq:dmkost}\boxed{Z=k\,\tau^{a}\quad\Longrightarrow\quad\tau_{\mathrm{req}}=\Bigl(\frac{Z_{\mathrm{req}}}{k}\Bigr)^{1/a}}\end{equation}

\begin{kfig}{The Kostiakov intake curve $Z=k\tau^a$: the infiltration rate is high at first and falls with time; the time to infiltrate the required depth is read from the curve.}
\draw[axis] (0,0)--(6.6,0);\draw[axis] (0,0)--(0,3.4);
\node[lab,below] at (3.2,-0.1) {opportunity time $\tau$};\node[lab,rotate=90] at (-0.4,1.6) {cumulative intake $Z$};
\draw[crv,smooth,domain=0:6.2] plot (\x,{1.25*\x^0.55});
\draw[gd] (0,2.6)--(3.78,2.6)--(3.78,0);\node[lab,left] at (0,2.6) {$Z_{\mathrm{req}}$};\node[lab,below] at (3.78,0) {$\tau_{\mathrm{req}}$};
\end{kfig}

\exbox{Time to infiltrate the required depth}{A soil has $k=0.012$ m\,min$^{-a}$ and $a=0.45$. How long must water stand on it to infiltrate 80 mm?}\sol
\[\tau_{\mathrm{req}}=\Bigl(\frac{0.080}{0.012}\Bigr)^{1/0.45}=6.667^{2.222}=67.7\ \text{min}\quad\text{\finans}\]
}{}
\section{Summary}
\begin{itemize}
\item ET$_0$ from temperature (Hargreaves), radiation (Makkink), or the service (Penman--Monteith).
\item A weather file is gridded by $\rho_{\rm ref}$ and $\Delta T_{\rm ref}$; the daily layers use the file as is.
\item Area and perimeter on the sphere; slope by a fitted plane; water source by an engineered score.
\item NDVI and companions from reflectances; zones from standardised series; $K_c$ from NDVI cover.
\item Runoff by curve number; gullies by the $SA^{0.4}$ threshold; soil loss by RUSLE.
\end{itemize}
\IfFileExists{extra/app06.tex}{\subsection{The output of the application, panel by panel}
The design mathematics of Chapter 6 rests on weather and field data; \emph{Reliability over 20 seasons} shows how a design that works in one year behaves in all of them.

For each of the 20 weather years the layer gives the irrigation, the drainage, the days of stress and the yield; a design is reliable if the yearly values stay within the band.

\appfig{yby_0}{0.96}{Year by year: how reliable is this setting? Yearly values of irrigation, drainage, evapotranspiration and yield.}

The bars show the spread of two quantities that govern the design: the yearly crop water use and the stress days, which are the basis of the percentiles used for the design discharge and the capacity.

\appfig{yby_1}{0.96}{Evapotranspiration (mm per year) and days of crop water stress per season, across the 20 years.}

The summary gives the share of the seasons in which the setting meets the target without stress, a direct reading of the reliability concept of the design chapter.

\appfig{yby_2}{0.96}{Reliability of this setting over the 20 simulated seasons from the yearly values above.}

}{}
\IfFileExists{extra/long06.tex}{\subsection{Complete calculation of a system capacity}
\exbox{Net and gross requirement, interval and capacity}{A 40 ha farm has a peak crop requirement of 6.2 mm\,d$^{-1}$, a net application of 45 mm, an application efficiency of 85 per cent and a leaching requirement of 10 per cent. Find the interval, the gross depth and volume, the system capacity for a 6-day cycle of 20 h days, and the seasonal volume.}\sol
\ans Peak crop need $ET_c=6.2$ mm\,d$^{-1}$; net application $Z_n=45$ mm (trigger of 50 per cent of 90 mm TAW)
\ans Interval $=Z_n/ET_c=45/6.2=7.3$ d; adopt 7 days
\ans Net depth every 7 days $=6.2\times7=43.4$ mm
\ans Gross depth with application efficiency 0.85 and a leaching fraction $LR=0.10$: $Z_g=Z_n/[E_a(1-LR)]=56.7$ mm
\ans Gross volume per irrigation for 40 ha: $10Z_gA=22693$ m$^3$
\ans System capacity needed to irrigate 40 ha in 6 days of 20 h: \[Q=V/(6\times20\times3600)\times1000=52.5\ \text{L\,s$^{-1}$}\]
\ans Seasonal depth over 150 days: $ET_c$ mean 5.0 mm\,d$^{-1}$ gives $5.0\times150/(0.85\times0.9)=980$ mm, or 392157 m$^3$ for 40 ha \finans


}{}
\IfFileExists{extra/proj06.tex}{\banner{Design project 6: sizing the mainline and the pump of a sprinkler block}
\textbf{Brief.} A sprinkler block needs $Q=60$ L\,s$^{-1}$ at 35 m of working head (operating pressure at the block plus lateral losses, and the lift). The mainline is 800 m of PVC ($C=140$). The designer chooses between standard diameters from 110 to 250 mm, balancing the cost of pipe (taken as $0.0004D^{1.7}$ currency units per metre, $D$ in mm, an assumed price law) against the cost of the energy to overcome the friction, with a pump efficiency of 65 per cent, 1\,200 h of pumping per year and an energy price of 0.12 units per kWh over a 15-year life at a discount of 6 per cent (present-value factor 9.71).

\begin{table}[!htbp]\centering\small\caption{Design table of the mainline: velocity, friction loss, pump power for friction plus 35 m, pipe cost, the check on velocity, present value of energy and total cost.}
\begin{tabular}{rrrrrlrr}\toprule
\hd{$D$ (mm)} & \hd{$v$ (m/s)} & \hd{$h_f$ (m)} & \hd{Power (kW)} & \hd{Pipe cost} & \hd{Velocity} & \hd{Energy PV} & \hd{Total}\\\midrule
110 & 6.31 & 230.7 & 240.6 & 945 & v>1.5 & 336417 & 337362\\
125 & 4.89 & 123.8 & 143.8 & 1175 & v>1.5 & 201067 & 202242\\
140 & 3.90 & 71.3 & 96.2 & 1424 & v>1.5 & 134511 & 135935\\
160 & 2.98 & 37.2 & 65.4 & 1787 & v>1.5 & 91445 & 93232\\
180 & 2.36 & 21.0 & 50.7 & 2183 & v>1.5 & 70891 & 73074\\
200 & 1.91 & 12.5 & 43.1 & 2612 & v>1.5 & 60264 & 62876\\
225 & 1.51 & 7.1 & 38.1 & 3191 & v>1.5 & 53273 & 56464\\
250 & 1.22 & 4.2 & 35.5 & 3816 & ok & 49638 & 53454\\\bottomrule\end{tabular}\end{table}

\textbf{Decision.} Among the diameters with an acceptable velocity the least total cost is for $D=250$ mm: velocity 1.22 m\,s$^{-1}$, friction loss 4.2 m, pump power 35.5 kW, total 53454 units. Smaller pipes are cheaper but the energy to push the water through them dominates; larger pipes cost more than the energy they save. The optimum is shallow, so a neighbouring standard size within a few per cent of the cost may be chosen for availability. The pump must deliver $60$ L\,s$^{-1}$ at a head of $35+4.2=39.2$ m.

\textbf{Check.} Water hammer: for a closing time of 10 s in a PVC pipe of wave speed 400 m\,s$^{-1}$ and length 800 m, the surge is $\Delta H=2Lv/(gt)=2\times800\times1.22/(9.81\times10)=20$ m, well above the 35 m working head: a slow-closing valve or an air vessel must be specified.

}{}
\IfFileExists{extra/ten06.tex}{\section{Worked solutions on the ten soils}
Whether a sprinkler rate produces runoff depends on the intake of the soil; Green--Ampt is used.

\exbox{Ten-soil example 6a: Coarse loam (S01)}{Coarse loam: a sprinkler system applies $i=14$ mm\,h$^{-1}$. With $K_s=13.0$ mm\,h$^{-1}$, $\psi_f=110$ mm and $\Delta\theta=0.34$ (class values), find whether ponding occurs, the time of ponding and the cumulative infiltration at that time (Green--Ampt).}\sol Since $i=14>K_s=13.0$, ponding occurs. $t_p=K_s\psi_f\Delta\theta/[i(i-K_s)]=13.0\times110\times0.34/[14\times1.0]=34.73$ h \ans $F_p=i\,t_p=14\times34.73=486.2$ mm \ans After ponding the surplus runs off or must be stored; the rate should be lowered or the set shortened.

\exbox{Ten-soil example 6a: Coarse sand (S02)}{Coarse sand: a sprinkler system applies $i=14$ mm\,h$^{-1}$. With $K_s=120.0$ mm\,h$^{-1}$, $\psi_f=49$ mm and $\Delta\theta=0.38$ (class values), find whether ponding occurs, the time of ponding and the cumulative infiltration at that time (Green--Ampt).}\sol Since $i=14\le K_s=120.0$ mm\,h$^{-1}$ the soil takes all the water at the surface: no ponding at any time \ans The limit is the nozzle, not the soil; the system can be run for any set time without runoff \ans

\exbox{Ten-soil example 6a: Sand (S03)}{Sand: a sprinkler system applies $i=14$ mm\,h$^{-1}$. With $K_s=210.0$ mm\,h$^{-1}$, $\psi_f=49$ mm and $\Delta\theta=0.38$ (class values), find whether ponding occurs, the time of ponding and the cumulative infiltration at that time (Green--Ampt).}\sol Since $i=14\le K_s=210.0$ mm\,h$^{-1}$ the soil takes all the water at the surface: no ponding at any time \ans The limit is the nozzle, not the soil; the system can be run for any set time without runoff \ans

\exbox{Ten-soil example 6a: Loamy sand (S04)}{Loamy sand: a sprinkler system applies $i=14$ mm\,h$^{-1}$. With $K_s=61.0$ mm\,h$^{-1}$, $\psi_f=61$ mm and $\Delta\theta=0.37$ (class values), find whether ponding occurs, the time of ponding and the cumulative infiltration at that time (Green--Ampt).}\sol Since $i=14\le K_s=61.0$ mm\,h$^{-1}$ the soil takes all the water at the surface: no ponding at any time \ans The limit is the nozzle, not the soil; the system can be run for any set time without runoff \ans

\exbox{Ten-soil example 6a: Sandy loam (S05)}{Sandy loam: a sprinkler system applies $i=14$ mm\,h$^{-1}$. With $K_s=26.0$ mm\,h$^{-1}$, $\psi_f=110$ mm and $\Delta\theta=0.35$ (class values), find whether ponding occurs, the time of ponding and the cumulative infiltration at that time (Green--Ampt).}\sol Since $i=14\le K_s=26.0$ mm\,h$^{-1}$ the soil takes all the water at the surface: no ponding at any time \ans The limit is the nozzle, not the soil; the system can be run for any set time without runoff \ans

\exbox{Ten-soil example 6a: Loam (S06)}{Loam: a sprinkler system applies $i=14$ mm\,h$^{-1}$. With $K_s=13.0$ mm\,h$^{-1}$, $\psi_f=89$ mm and $\Delta\theta=0.32$ (class values), find whether ponding occurs, the time of ponding and the cumulative infiltration at that time (Green--Ampt).}\sol Since $i=14>K_s=13.0$, ponding occurs. $t_p=K_s\psi_f\Delta\theta/[i(i-K_s)]=13.0\times89\times0.32/[14\times1.0]=26.45$ h \ans $F_p=i\,t_p=14\times26.45=370.2$ mm \ans After ponding the surplus runs off or must be stored; the rate should be lowered or the set shortened.

\exbox{Ten-soil example 6a: Silt loam (S07)}{Silt loam: a sprinkler system applies $i=14$ mm\,h$^{-1}$. With $K_s=6.8$ mm\,h$^{-1}$, $\psi_f=170$ mm and $\Delta\theta=0.35$ (class values), find whether ponding occurs, the time of ponding and the cumulative infiltration at that time (Green--Ampt).}\sol Since $i=14>K_s=6.8$, ponding occurs. $t_p=K_s\psi_f\Delta\theta/[i(i-K_s)]=6.8\times170\times0.35/[14\times7.2]=4.01$ h \ans $F_p=i\,t_p=14\times4.01=56.2$ mm \ans After ponding the surplus runs off or must be stored; the rate should be lowered or the set shortened.

\exbox{Ten-soil example 6a: Sandy clay loam (S08)}{Sandy clay loam: a sprinkler system applies $i=14$ mm\,h$^{-1}$. With $K_s=4.3$ mm\,h$^{-1}$, $\psi_f=219$ mm and $\Delta\theta=0.30$ (class values), find whether ponding occurs, the time of ponding and the cumulative infiltration at that time (Green--Ampt).}\sol Since $i=14>K_s=4.3$, ponding occurs. $t_p=K_s\psi_f\Delta\theta/[i(i-K_s)]=4.3\times219\times0.30/[14\times9.7]=2.08$ h \ans $F_p=i\,t_p=14\times2.08=29.1$ mm \ans After ponding the surplus runs off or must be stored; the rate should be lowered or the set shortened.

\exbox{Ten-soil example 6a: Clay loam (S09)}{Clay loam: a sprinkler system applies $i=14$ mm\,h$^{-1}$. With $K_s=2.3$ mm\,h$^{-1}$, $\psi_f=209$ mm and $\Delta\theta=0.31$ (class values), find whether ponding occurs, the time of ponding and the cumulative infiltration at that time (Green--Ampt).}\sol Since $i=14>K_s=2.3$, ponding occurs. $t_p=K_s\psi_f\Delta\theta/[i(i-K_s)]=2.3\times209\times0.31/[14\times11.7]=0.91$ h \ans $F_p=i\,t_p=14\times0.91=12.7$ mm \ans After ponding the surplus runs off or must be stored; the rate should be lowered or the set shortened.

\exbox{Ten-soil example 6a: Silty clay loam (S10)}{Silty clay loam: a sprinkler system applies $i=14$ mm\,h$^{-1}$. With $K_s=1.5$ mm\,h$^{-1}$, $\psi_f=273$ mm and $\Delta\theta=0.30$ (class values), find whether ponding occurs, the time of ponding and the cumulative infiltration at that time (Green--Ampt).}\sol Since $i=14>K_s=1.5$, ponding occurs. $t_p=K_s\psi_f\Delta\theta/[i(i-K_s)]=1.5\times273\times0.30/[14\times12.5]=0.70$ h \ans $F_p=i\,t_p=14\times0.70=9.8$ mm \ans After ponding the surplus runs off or must be stored; the rate should be lowered or the set shortened.

}{}
\IfFileExists{extra/tenb06.tex}{\section{Worked solutions on ten crops}
The ten crops of the library differ in their demand, their roots and their sensitivity. The same type of question is answered for each of them at $ET_0=5.5$ mm\,d$^{-1}$; the data are typical values of the crops (Doorenbos--Kassam and FAO 56).

\exbox{Crop example 6a: Maize}{Maize: find the peak crop evapotranspiration, the gross depth per day and the continuous flow per hectare for $E_a=0.75$ and $ET_0=5.5$ mm\,d$^{-1}$.}\sol $ET_c=1.20\times5.5=6.60$ mm\,d$^{-1}$ \ans Gross $=6.60/0.75=8.80$ mm\,d$^{-1}$ \ans Flow $=10\times8.80/24\,\mathrm{m^3\,h^{-1}\,ha^{-1}}=3.67$ m$^3$\,h$^{-1}$, i.e.\ $1.02$ L\,s$^{-1}$\,ha$^{-1}$ \ans

\exbox{Crop example 6a: Wheat}{Wheat: find the peak crop evapotranspiration, the gross depth per day and the continuous flow per hectare for $E_a=0.75$ and $ET_0=5.5$ mm\,d$^{-1}$.}\sol $ET_c=1.15\times5.5=6.32$ mm\,d$^{-1}$ \ans Gross $=6.32/0.75=8.43$ mm\,d$^{-1}$ \ans Flow $=10\times8.43/24\,\mathrm{m^3\,h^{-1}\,ha^{-1}}=3.51$ m$^3$\,h$^{-1}$, i.e.\ $0.98$ L\,s$^{-1}$\,ha$^{-1}$ \ans

\exbox{Crop example 6a: Cotton}{Cotton: find the peak crop evapotranspiration, the gross depth per day and the continuous flow per hectare for $E_a=0.75$ and $ET_0=5.5$ mm\,d$^{-1}$.}\sol $ET_c=1.15\times5.5=6.32$ mm\,d$^{-1}$ \ans Gross $=6.32/0.75=8.43$ mm\,d$^{-1}$ \ans Flow $=10\times8.43/24\,\mathrm{m^3\,h^{-1}\,ha^{-1}}=3.51$ m$^3$\,h$^{-1}$, i.e.\ $0.98$ L\,s$^{-1}$\,ha$^{-1}$ \ans

\exbox{Crop example 6a: Tomato}{Tomato: find the peak crop evapotranspiration, the gross depth per day and the continuous flow per hectare for $E_a=0.75$ and $ET_0=5.5$ mm\,d$^{-1}$.}\sol $ET_c=1.15\times5.5=6.32$ mm\,d$^{-1}$ \ans Gross $=6.32/0.75=8.43$ mm\,d$^{-1}$ \ans Flow $=10\times8.43/24\,\mathrm{m^3\,h^{-1}\,ha^{-1}}=3.51$ m$^3$\,h$^{-1}$, i.e.\ $0.98$ L\,s$^{-1}$\,ha$^{-1}$ \ans

\exbox{Crop example 6a: Potato}{Potato: find the peak crop evapotranspiration, the gross depth per day and the continuous flow per hectare for $E_a=0.75$ and $ET_0=5.5$ mm\,d$^{-1}$.}\sol $ET_c=1.15\times5.5=6.32$ mm\,d$^{-1}$ \ans Gross $=6.32/0.75=8.43$ mm\,d$^{-1}$ \ans Flow $=10\times8.43/24\,\mathrm{m^3\,h^{-1}\,ha^{-1}}=3.51$ m$^3$\,h$^{-1}$, i.e.\ $0.98$ L\,s$^{-1}$\,ha$^{-1}$ \ans

\exbox{Crop example 6a: Soybean}{Soybean: find the peak crop evapotranspiration, the gross depth per day and the continuous flow per hectare for $E_a=0.75$ and $ET_0=5.5$ mm\,d$^{-1}$.}\sol $ET_c=1.15\times5.5=6.32$ mm\,d$^{-1}$ \ans Gross $=6.32/0.75=8.43$ mm\,d$^{-1}$ \ans Flow $=10\times8.43/24\,\mathrm{m^3\,h^{-1}\,ha^{-1}}=3.51$ m$^3$\,h$^{-1}$, i.e.\ $0.98$ L\,s$^{-1}$\,ha$^{-1}$ \ans

\exbox{Crop example 6a: Sugar beet}{Sugar beet: find the peak crop evapotranspiration, the gross depth per day and the continuous flow per hectare for $E_a=0.75$ and $ET_0=5.5$ mm\,d$^{-1}$.}\sol $ET_c=1.20\times5.5=6.60$ mm\,d$^{-1}$ \ans Gross $=6.60/0.75=8.80$ mm\,d$^{-1}$ \ans Flow $=10\times8.80/24\,\mathrm{m^3\,h^{-1}\,ha^{-1}}=3.67$ m$^3$\,h$^{-1}$, i.e.\ $1.02$ L\,s$^{-1}$\,ha$^{-1}$ \ans

\exbox{Crop example 6a: Sunflower}{Sunflower: find the peak crop evapotranspiration, the gross depth per day and the continuous flow per hectare for $E_a=0.75$ and $ET_0=5.5$ mm\,d$^{-1}$.}\sol $ET_c=1.10\times5.5=6.05$ mm\,d$^{-1}$ \ans Gross $=6.05/0.75=8.07$ mm\,d$^{-1}$ \ans Flow $=10\times8.07/24\,\mathrm{m^3\,h^{-1}\,ha^{-1}}=3.36$ m$^3$\,h$^{-1}$, i.e.\ $0.93$ L\,s$^{-1}$\,ha$^{-1}$ \ans

\exbox{Crop example 6a: Onion}{Onion: find the peak crop evapotranspiration, the gross depth per day and the continuous flow per hectare for $E_a=0.75$ and $ET_0=5.5$ mm\,d$^{-1}$.}\sol $ET_c=1.05\times5.5=5.78$ mm\,d$^{-1}$ \ans Gross $=5.78/0.75=7.70$ mm\,d$^{-1}$ \ans Flow $=10\times7.70/24\,\mathrm{m^3\,h^{-1}\,ha^{-1}}=3.21$ m$^3$\,h$^{-1}$, i.e.\ $0.89$ L\,s$^{-1}$\,ha$^{-1}$ \ans

\exbox{Crop example 6a: Alfalfa}{Alfalfa: find the peak crop evapotranspiration, the gross depth per day and the continuous flow per hectare for $E_a=0.75$ and $ET_0=5.5$ mm\,d$^{-1}$.}\sol $ET_c=1.15\times5.5=6.32$ mm\,d$^{-1}$ \ans Gross $=6.32/0.75=8.43$ mm\,d$^{-1}$ \ans Flow $=10\times8.43/24\,\mathrm{m^3\,h^{-1}\,ha^{-1}}=3.51$ m$^3$\,h$^{-1}$, i.e.\ $0.98$ L\,s$^{-1}$\,ha$^{-1}$ \ans

}{}
\IfFileExists{extra/teny06.tex}{\subsection{Worked solutions on the ten soils (third scenario)}
The same questions are repeated for different working conditions, to show how the answer moves with them.

\exbox{Ten-soil example 6c: Coarse loam (S01)}{Coarse loam: a sprinkler system applies $i=20$ mm\,h$^{-1}$. With $K_s=13.0$ mm\,h$^{-1}$, $\psi_f=110$ mm and $\Delta\theta=0.34$ (class values), find whether ponding occurs, the time of ponding and the cumulative infiltration at that time (Green--Ampt).}\sol \givenline{$i=20$ mm\,h$^{-1}$ ; $K_s=13.0$ mm\,h$^{-1}$ ; $\psi_f=110$ mm ; $\Delta\theta=0.34$}\ans Since $i=20>K_s=13.0$, ponding occurs. $t_p=K_s\psi_f\Delta\theta/[i(i-K_s)]=13.0\times110\times0.34/[20\times7.0]=3.47$ h \ans $F_p=i\,t_p=20\times3.47=69.5$ mm \ans After ponding the surplus runs off or must be stored; the rate should be lowered or the set shortened. \finans

\exbox{Ten-soil example 6c: Coarse sand (S02)}{Coarse sand: a sprinkler system applies $i=20$ mm\,h$^{-1}$. With $K_s=120.0$ mm\,h$^{-1}$, $\psi_f=49$ mm and $\Delta\theta=0.38$ (class values), find whether ponding occurs, the time of ponding and the cumulative infiltration at that time (Green--Ampt).}\sol \givenline{$i=20$ mm\,h$^{-1}$ ; $K_s=120.0$ mm\,h$^{-1}$ ; $\psi_f=49$ mm ; $\Delta\theta=0.38$}\ans Since $i=20\le K_s=120.0$ mm\,h$^{-1}$ the soil takes all the water at the surface: no ponding at any time \ans The limit is the nozzle, not the soil; the system can be run for any set time without runoff \finans

\exbox{Ten-soil example 6c: Sand (S03)}{Sand: a sprinkler system applies $i=20$ mm\,h$^{-1}$. With $K_s=210.0$ mm\,h$^{-1}$, $\psi_f=49$ mm and $\Delta\theta=0.38$ (class values), find whether ponding occurs, the time of ponding and the cumulative infiltration at that time (Green--Ampt).}\sol \givenline{$i=20$ mm\,h$^{-1}$ ; $K_s=210.0$ mm\,h$^{-1}$ ; $\psi_f=49$ mm ; $\Delta\theta=0.38$}\ans Since $i=20\le K_s=210.0$ mm\,h$^{-1}$ the soil takes all the water at the surface: no ponding at any time \ans The limit is the nozzle, not the soil; the system can be run for any set time without runoff \finans

\exbox{Ten-soil example 6c: Loamy sand (S04)}{Loamy sand: a sprinkler system applies $i=20$ mm\,h$^{-1}$. With $K_s=61.0$ mm\,h$^{-1}$, $\psi_f=61$ mm and $\Delta\theta=0.37$ (class values), find whether ponding occurs, the time of ponding and the cumulative infiltration at that time (Green--Ampt).}\sol \givenline{$i=20$ mm\,h$^{-1}$ ; $K_s=61.0$ mm\,h$^{-1}$ ; $\psi_f=61$ mm ; $\Delta\theta=0.37$}\ans Since $i=20\le K_s=61.0$ mm\,h$^{-1}$ the soil takes all the water at the surface: no ponding at any time \ans The limit is the nozzle, not the soil; the system can be run for any set time without runoff \finans

\exbox{Ten-soil example 6c: Sandy loam (S05)}{Sandy loam: a sprinkler system applies $i=20$ mm\,h$^{-1}$. With $K_s=26.0$ mm\,h$^{-1}$, $\psi_f=110$ mm and $\Delta\theta=0.35$ (class values), find whether ponding occurs, the time of ponding and the cumulative infiltration at that time (Green--Ampt).}\sol \givenline{$i=20$ mm\,h$^{-1}$ ; $K_s=26.0$ mm\,h$^{-1}$ ; $\psi_f=110$ mm ; $\Delta\theta=0.35$}\ans Since $i=20\le K_s=26.0$ mm\,h$^{-1}$ the soil takes all the water at the surface: no ponding at any time \ans The limit is the nozzle, not the soil; the system can be run for any set time without runoff \finans

\exbox{Ten-soil example 6c: Loam (S06)}{Loam: a sprinkler system applies $i=20$ mm\,h$^{-1}$. With $K_s=13.0$ mm\,h$^{-1}$, $\psi_f=89$ mm and $\Delta\theta=0.32$ (class values), find whether ponding occurs, the time of ponding and the cumulative infiltration at that time (Green--Ampt).}\sol \givenline{$i=20$ mm\,h$^{-1}$ ; $K_s=13.0$ mm\,h$^{-1}$ ; $\psi_f=89$ mm ; $\Delta\theta=0.32$}\ans Since $i=20>K_s=13.0$, ponding occurs. $t_p=K_s\psi_f\Delta\theta/[i(i-K_s)]=13.0\times89\times0.32/[20\times7.0]=2.64$ h \ans $F_p=i\,t_p=20\times2.64=52.9$ mm \ans After ponding the surplus runs off or must be stored; the rate should be lowered or the set shortened. \finans

\exbox{Ten-soil example 6c: Silt loam (S07)}{Silt loam: a sprinkler system applies $i=20$ mm\,h$^{-1}$. With $K_s=6.8$ mm\,h$^{-1}$, $\psi_f=170$ mm and $\Delta\theta=0.35$ (class values), find whether ponding occurs, the time of ponding and the cumulative infiltration at that time (Green--Ampt).}\sol \givenline{$i=20$ mm\,h$^{-1}$ ; $K_s=6.8$ mm\,h$^{-1}$ ; $\psi_f=170$ mm ; $\Delta\theta=0.35$}\ans Since $i=20>K_s=6.8$, ponding occurs. $t_p=K_s\psi_f\Delta\theta/[i(i-K_s)]=6.8\times170\times0.35/[20\times13.2]=1.53$ h \ans $F_p=i\,t_p=20\times1.53=30.7$ mm \ans After ponding the surplus runs off or must be stored; the rate should be lowered or the set shortened. \finans

\exbox{Ten-soil example 6c: Sandy clay loam (S08)}{Sandy clay loam: a sprinkler system applies $i=20$ mm\,h$^{-1}$. With $K_s=4.3$ mm\,h$^{-1}$, $\psi_f=219$ mm and $\Delta\theta=0.30$ (class values), find whether ponding occurs, the time of ponding and the cumulative infiltration at that time (Green--Ampt).}\sol \givenline{$i=20$ mm\,h$^{-1}$ ; $K_s=4.3$ mm\,h$^{-1}$ ; $\psi_f=219$ mm ; $\Delta\theta=0.30$}\ans Since $i=20>K_s=4.3$, ponding occurs. $t_p=K_s\psi_f\Delta\theta/[i(i-K_s)]=4.3\times219\times0.30/[20\times15.7]=0.90$ h \ans $F_p=i\,t_p=20\times0.90=18.0$ mm \ans After ponding the surplus runs off or must be stored; the rate should be lowered or the set shortened. \finans

\exbox{Ten-soil example 6c: Clay loam (S09)}{Clay loam: a sprinkler system applies $i=20$ mm\,h$^{-1}$. With $K_s=2.3$ mm\,h$^{-1}$, $\psi_f=209$ mm and $\Delta\theta=0.31$ (class values), find whether ponding occurs, the time of ponding and the cumulative infiltration at that time (Green--Ampt).}\sol \givenline{$i=20$ mm\,h$^{-1}$ ; $K_s=2.3$ mm\,h$^{-1}$ ; $\psi_f=209$ mm ; $\Delta\theta=0.31$}\ans Since $i=20>K_s=2.3$, ponding occurs. $t_p=K_s\psi_f\Delta\theta/[i(i-K_s)]=2.3\times209\times0.31/[20\times17.7]=0.42$ h \ans $F_p=i\,t_p=20\times0.42=8.4$ mm \ans After ponding the surplus runs off or must be stored; the rate should be lowered or the set shortened. \finans

\exbox{Ten-soil example 6c: Silty clay loam (S10)}{Silty clay loam: a sprinkler system applies $i=20$ mm\,h$^{-1}$. With $K_s=1.5$ mm\,h$^{-1}$, $\psi_f=273$ mm and $\Delta\theta=0.30$ (class values), find whether ponding occurs, the time of ponding and the cumulative infiltration at that time (Green--Ampt).}\sol \givenline{$i=20$ mm\,h$^{-1}$ ; $K_s=1.5$ mm\,h$^{-1}$ ; $\psi_f=273$ mm ; $\Delta\theta=0.30$}\ans Since $i=20>K_s=1.5$, ponding occurs. $t_p=K_s\psi_f\Delta\theta/[i(i-K_s)]=1.5\times273\times0.30/[20\times18.5]=0.33$ h \ans $F_p=i\,t_p=20\times0.33=6.6$ mm \ans After ponding the surplus runs off or must be stored; the rate should be lowered or the set shortened. \finans

}{}
\IfFileExists{extra/tenby06.tex}{\subsection{Worked solutions on ten crops (hot climate)}
The ten crops of the library differ in their demand, their roots and their sensitivity. The same type of question is answered for each of them at $ET_0=7.5$ mm\,d$^{-1}$; the data are typical values of the crops (Doorenbos--Kassam and FAO 56).

\exbox{Crop example 6c: Maize}{Maize: find the peak crop evapotranspiration, the gross depth per day and the continuous flow per hectare for $E_a=0.75$ and $ET_0=7.5$ mm\,d$^{-1}$.}\sol \givenline{$E_a=0.75$ ; $ET_0=7.5$ mm\,d$^{-1}$}\ans $ET_c=1.20\times7.5=9.00$ mm\,d$^{-1}$ \ans Gross $=9.00/0.75=12.00$ mm\,d$^{-1}$ \ans Flow $=10\times12.00/24\,\mathrm{m^3\,h^{-1}\,ha^{-1}}=5.00$ m$^3$\,h$^{-1}$, i.e.\ $1.39$ L\,s$^{-1}$\,ha$^{-1}$ \finans

\exbox{Crop example 6c: Wheat}{Wheat: find the peak crop evapotranspiration, the gross depth per day and the continuous flow per hectare for $E_a=0.75$ and $ET_0=7.5$ mm\,d$^{-1}$.}\sol \givenline{$E_a=0.75$ ; $ET_0=7.5$ mm\,d$^{-1}$}\ans $ET_c=1.15\times7.5=8.62$ mm\,d$^{-1}$ \ans Gross $=8.62/0.75=11.50$ mm\,d$^{-1}$ \ans Flow $=10\times11.50/24\,\mathrm{m^3\,h^{-1}\,ha^{-1}}=4.79$ m$^3$\,h$^{-1}$, i.e.\ $1.33$ L\,s$^{-1}$\,ha$^{-1}$ \finans

\exbox{Crop example 6c: Cotton}{Cotton: find the peak crop evapotranspiration, the gross depth per day and the continuous flow per hectare for $E_a=0.75$ and $ET_0=7.5$ mm\,d$^{-1}$.}\sol \givenline{$E_a=0.75$ ; $ET_0=7.5$ mm\,d$^{-1}$}\ans $ET_c=1.15\times7.5=8.62$ mm\,d$^{-1}$ \ans Gross $=8.62/0.75=11.50$ mm\,d$^{-1}$ \ans Flow $=10\times11.50/24\,\mathrm{m^3\,h^{-1}\,ha^{-1}}=4.79$ m$^3$\,h$^{-1}$, i.e.\ $1.33$ L\,s$^{-1}$\,ha$^{-1}$ \finans

\exbox{Crop example 6c: Tomato}{Tomato: find the peak crop evapotranspiration, the gross depth per day and the continuous flow per hectare for $E_a=0.75$ and $ET_0=7.5$ mm\,d$^{-1}$.}\sol \givenline{$E_a=0.75$ ; $ET_0=7.5$ mm\,d$^{-1}$}\ans $ET_c=1.15\times7.5=8.62$ mm\,d$^{-1}$ \ans Gross $=8.62/0.75=11.50$ mm\,d$^{-1}$ \ans Flow $=10\times11.50/24\,\mathrm{m^3\,h^{-1}\,ha^{-1}}=4.79$ m$^3$\,h$^{-1}$, i.e.\ $1.33$ L\,s$^{-1}$\,ha$^{-1}$ \finans

\exbox{Crop example 6c: Potato}{Potato: find the peak crop evapotranspiration, the gross depth per day and the continuous flow per hectare for $E_a=0.75$ and $ET_0=7.5$ mm\,d$^{-1}$.}\sol \givenline{$E_a=0.75$ ; $ET_0=7.5$ mm\,d$^{-1}$}\ans $ET_c=1.15\times7.5=8.62$ mm\,d$^{-1}$ \ans Gross $=8.62/0.75=11.50$ mm\,d$^{-1}$ \ans Flow $=10\times11.50/24\,\mathrm{m^3\,h^{-1}\,ha^{-1}}=4.79$ m$^3$\,h$^{-1}$, i.e.\ $1.33$ L\,s$^{-1}$\,ha$^{-1}$ \finans

\exbox{Crop example 6c: Soybean}{Soybean: find the peak crop evapotranspiration, the gross depth per day and the continuous flow per hectare for $E_a=0.75$ and $ET_0=7.5$ mm\,d$^{-1}$.}\sol \givenline{$E_a=0.75$ ; $ET_0=7.5$ mm\,d$^{-1}$}\ans $ET_c=1.15\times7.5=8.62$ mm\,d$^{-1}$ \ans Gross $=8.62/0.75=11.50$ mm\,d$^{-1}$ \ans Flow $=10\times11.50/24\,\mathrm{m^3\,h^{-1}\,ha^{-1}}=4.79$ m$^3$\,h$^{-1}$, i.e.\ $1.33$ L\,s$^{-1}$\,ha$^{-1}$ \finans

\exbox{Crop example 6c: Sugar beet}{Sugar beet: find the peak crop evapotranspiration, the gross depth per day and the continuous flow per hectare for $E_a=0.75$ and $ET_0=7.5$ mm\,d$^{-1}$.}\sol \givenline{$E_a=0.75$ ; $ET_0=7.5$ mm\,d$^{-1}$}\ans $ET_c=1.20\times7.5=9.00$ mm\,d$^{-1}$ \ans Gross $=9.00/0.75=12.00$ mm\,d$^{-1}$ \ans Flow $=10\times12.00/24\,\mathrm{m^3\,h^{-1}\,ha^{-1}}=5.00$ m$^3$\,h$^{-1}$, i.e.\ $1.39$ L\,s$^{-1}$\,ha$^{-1}$ \finans

\exbox{Crop example 6c: Sunflower}{Sunflower: find the peak crop evapotranspiration, the gross depth per day and the continuous flow per hectare for $E_a=0.75$ and $ET_0=7.5$ mm\,d$^{-1}$.}\sol \givenline{$E_a=0.75$ ; $ET_0=7.5$ mm\,d$^{-1}$}\ans $ET_c=1.10\times7.5=8.25$ mm\,d$^{-1}$ \ans Gross $=8.25/0.75=11.00$ mm\,d$^{-1}$ \ans Flow $=10\times11.00/24\,\mathrm{m^3\,h^{-1}\,ha^{-1}}=4.58$ m$^3$\,h$^{-1}$, i.e.\ $1.27$ L\,s$^{-1}$\,ha$^{-1}$ \finans

\exbox{Crop example 6c: Onion}{Onion: find the peak crop evapotranspiration, the gross depth per day and the continuous flow per hectare for $E_a=0.75$ and $ET_0=7.5$ mm\,d$^{-1}$.}\sol \givenline{$E_a=0.75$ ; $ET_0=7.5$ mm\,d$^{-1}$}\ans $ET_c=1.05\times7.5=7.88$ mm\,d$^{-1}$ \ans Gross $=7.88/0.75=10.50$ mm\,d$^{-1}$ \ans Flow $=10\times10.50/24\,\mathrm{m^3\,h^{-1}\,ha^{-1}}=4.38$ m$^3$\,h$^{-1}$, i.e.\ $1.22$ L\,s$^{-1}$\,ha$^{-1}$ \finans

\exbox{Crop example 6c: Alfalfa}{Alfalfa: find the peak crop evapotranspiration, the gross depth per day and the continuous flow per hectare for $E_a=0.75$ and $ET_0=7.5$ mm\,d$^{-1}$.}\sol \givenline{$E_a=0.75$ ; $ET_0=7.5$ mm\,d$^{-1}$}\ans $ET_c=1.15\times7.5=8.62$ mm\,d$^{-1}$ \ans Gross $=8.62/0.75=11.50$ mm\,d$^{-1}$ \ans Flow $=10\times11.50/24\,\mathrm{m^3\,h^{-1}\,ha^{-1}}=4.79$ m$^3$\,h$^{-1}$, i.e.\ $1.33$ L\,s$^{-1}$\,ha$^{-1}$ \finans

}{}
\IfFileExists{extra/tenx06.tex}{\subsection{Worked solutions on the ten soils (second scenario)}
The same questions are repeated for different working conditions, to show how the answer moves with them.

\exbox{Ten-soil example 6b: Coarse loam (S01)}{Coarse loam: a sprinkler system applies $i=10$ mm\,h$^{-1}$. With $K_s=13.0$ mm\,h$^{-1}$, $\psi_f=110$ mm and $\Delta\theta=0.34$ (class values), find whether ponding occurs, the time of ponding and the cumulative infiltration at that time (Green--Ampt).}\sol Since $i=10\le K_s=13.0$ mm\,h$^{-1}$ the soil takes all the water at the surface: no ponding at any time \ans The limit is the nozzle, not the soil; the system can be run for any set time without runoff \ans

\exbox{Ten-soil example 6b: Coarse sand (S02)}{Coarse sand: a sprinkler system applies $i=10$ mm\,h$^{-1}$. With $K_s=120.0$ mm\,h$^{-1}$, $\psi_f=49$ mm and $\Delta\theta=0.38$ (class values), find whether ponding occurs, the time of ponding and the cumulative infiltration at that time (Green--Ampt).}\sol Since $i=10\le K_s=120.0$ mm\,h$^{-1}$ the soil takes all the water at the surface: no ponding at any time \ans The limit is the nozzle, not the soil; the system can be run for any set time without runoff \ans

\exbox{Ten-soil example 6b: Sand (S03)}{Sand: a sprinkler system applies $i=10$ mm\,h$^{-1}$. With $K_s=210.0$ mm\,h$^{-1}$, $\psi_f=49$ mm and $\Delta\theta=0.38$ (class values), find whether ponding occurs, the time of ponding and the cumulative infiltration at that time (Green--Ampt).}\sol Since $i=10\le K_s=210.0$ mm\,h$^{-1}$ the soil takes all the water at the surface: no ponding at any time \ans The limit is the nozzle, not the soil; the system can be run for any set time without runoff \ans

\exbox{Ten-soil example 6b: Loamy sand (S04)}{Loamy sand: a sprinkler system applies $i=10$ mm\,h$^{-1}$. With $K_s=61.0$ mm\,h$^{-1}$, $\psi_f=61$ mm and $\Delta\theta=0.37$ (class values), find whether ponding occurs, the time of ponding and the cumulative infiltration at that time (Green--Ampt).}\sol Since $i=10\le K_s=61.0$ mm\,h$^{-1}$ the soil takes all the water at the surface: no ponding at any time \ans The limit is the nozzle, not the soil; the system can be run for any set time without runoff \ans

\exbox{Ten-soil example 6b: Sandy loam (S05)}{Sandy loam: a sprinkler system applies $i=10$ mm\,h$^{-1}$. With $K_s=26.0$ mm\,h$^{-1}$, $\psi_f=110$ mm and $\Delta\theta=0.35$ (class values), find whether ponding occurs, the time of ponding and the cumulative infiltration at that time (Green--Ampt).}\sol Since $i=10\le K_s=26.0$ mm\,h$^{-1}$ the soil takes all the water at the surface: no ponding at any time \ans The limit is the nozzle, not the soil; the system can be run for any set time without runoff \ans

\exbox{Ten-soil example 6b: Loam (S06)}{Loam: a sprinkler system applies $i=10$ mm\,h$^{-1}$. With $K_s=13.0$ mm\,h$^{-1}$, $\psi_f=89$ mm and $\Delta\theta=0.32$ (class values), find whether ponding occurs, the time of ponding and the cumulative infiltration at that time (Green--Ampt).}\sol Since $i=10\le K_s=13.0$ mm\,h$^{-1}$ the soil takes all the water at the surface: no ponding at any time \ans The limit is the nozzle, not the soil; the system can be run for any set time without runoff \ans

\exbox{Ten-soil example 6b: Silt loam (S07)}{Silt loam: a sprinkler system applies $i=10$ mm\,h$^{-1}$. With $K_s=6.8$ mm\,h$^{-1}$, $\psi_f=170$ mm and $\Delta\theta=0.35$ (class values), find whether ponding occurs, the time of ponding and the cumulative infiltration at that time (Green--Ampt).}\sol Since $i=10>K_s=6.8$, ponding occurs. $t_p=K_s\psi_f\Delta\theta/[i(i-K_s)]=6.8\times170\times0.35/[10\times3.2]=12.64$ h \ans $F_p=i\,t_p=10\times12.64=126.4$ mm \ans After ponding the surplus runs off or must be stored; the rate should be lowered or the set shortened.

\exbox{Ten-soil example 6b: Sandy clay loam (S08)}{Sandy clay loam: a sprinkler system applies $i=10$ mm\,h$^{-1}$. With $K_s=4.3$ mm\,h$^{-1}$, $\psi_f=219$ mm and $\Delta\theta=0.30$ (class values), find whether ponding occurs, the time of ponding and the cumulative infiltration at that time (Green--Ampt).}\sol Since $i=10>K_s=4.3$, ponding occurs. $t_p=K_s\psi_f\Delta\theta/[i(i-K_s)]=4.3\times219\times0.30/[10\times5.7]=4.96$ h \ans $F_p=i\,t_p=10\times4.96=49.6$ mm \ans After ponding the surplus runs off or must be stored; the rate should be lowered or the set shortened.

\exbox{Ten-soil example 6b: Clay loam (S09)}{Clay loam: a sprinkler system applies $i=10$ mm\,h$^{-1}$. With $K_s=2.3$ mm\,h$^{-1}$, $\psi_f=209$ mm and $\Delta\theta=0.31$ (class values), find whether ponding occurs, the time of ponding and the cumulative infiltration at that time (Green--Ampt).}\sol Since $i=10>K_s=2.3$, ponding occurs. $t_p=K_s\psi_f\Delta\theta/[i(i-K_s)]=2.3\times209\times0.31/[10\times7.7]=1.94$ h \ans $F_p=i\,t_p=10\times1.94=19.4$ mm \ans After ponding the surplus runs off or must be stored; the rate should be lowered or the set shortened.

\exbox{Ten-soil example 6b: Silty clay loam (S10)}{Silty clay loam: a sprinkler system applies $i=10$ mm\,h$^{-1}$. With $K_s=1.5$ mm\,h$^{-1}$, $\psi_f=273$ mm and $\Delta\theta=0.30$ (class values), find whether ponding occurs, the time of ponding and the cumulative infiltration at that time (Green--Ampt).}\sol Since $i=10>K_s=1.5$, ponding occurs. $t_p=K_s\psi_f\Delta\theta/[i(i-K_s)]=1.5\times273\times0.30/[10\times8.5]=1.45$ h \ans $F_p=i\,t_p=10\times1.45=14.5$ mm \ans After ponding the surplus runs off or must be stored; the rate should be lowered or the set shortened.

}{}
\IfFileExists{extra/tenbx06.tex}{\subsection{Worked solutions on ten crops (cool climate)}
The ten crops of the library differ in their demand, their roots and their sensitivity. The same type of question is answered for each of them at $ET_0=4.0$ mm\,d$^{-1}$; the data are typical values of the crops (Doorenbos--Kassam and FAO 56).

\exbox{Crop example 6b: Maize}{Maize: find the peak crop evapotranspiration, the gross depth per day and the continuous flow per hectare for $E_a=0.75$ and $ET_0=4.0$ mm\,d$^{-1}$.}\sol $ET_c=1.20\times4.0=4.80$ mm\,d$^{-1}$ \ans Gross $=4.80/0.75=6.40$ mm\,d$^{-1}$ \ans Flow $=10\times6.40/24\,\mathrm{m^3\,h^{-1}\,ha^{-1}}=2.67$ m$^3$\,h$^{-1}$, i.e.\ $0.74$ L\,s$^{-1}$\,ha$^{-1}$ \ans

\exbox{Crop example 6b: Wheat}{Wheat: find the peak crop evapotranspiration, the gross depth per day and the continuous flow per hectare for $E_a=0.75$ and $ET_0=4.0$ mm\,d$^{-1}$.}\sol $ET_c=1.15\times4.0=4.60$ mm\,d$^{-1}$ \ans Gross $=4.60/0.75=6.13$ mm\,d$^{-1}$ \ans Flow $=10\times6.13/24\,\mathrm{m^3\,h^{-1}\,ha^{-1}}=2.56$ m$^3$\,h$^{-1}$, i.e.\ $0.71$ L\,s$^{-1}$\,ha$^{-1}$ \ans

\exbox{Crop example 6b: Cotton}{Cotton: find the peak crop evapotranspiration, the gross depth per day and the continuous flow per hectare for $E_a=0.75$ and $ET_0=4.0$ mm\,d$^{-1}$.}\sol $ET_c=1.15\times4.0=4.60$ mm\,d$^{-1}$ \ans Gross $=4.60/0.75=6.13$ mm\,d$^{-1}$ \ans Flow $=10\times6.13/24\,\mathrm{m^3\,h^{-1}\,ha^{-1}}=2.56$ m$^3$\,h$^{-1}$, i.e.\ $0.71$ L\,s$^{-1}$\,ha$^{-1}$ \ans

\exbox{Crop example 6b: Tomato}{Tomato: find the peak crop evapotranspiration, the gross depth per day and the continuous flow per hectare for $E_a=0.75$ and $ET_0=4.0$ mm\,d$^{-1}$.}\sol $ET_c=1.15\times4.0=4.60$ mm\,d$^{-1}$ \ans Gross $=4.60/0.75=6.13$ mm\,d$^{-1}$ \ans Flow $=10\times6.13/24\,\mathrm{m^3\,h^{-1}\,ha^{-1}}=2.56$ m$^3$\,h$^{-1}$, i.e.\ $0.71$ L\,s$^{-1}$\,ha$^{-1}$ \ans

\exbox{Crop example 6b: Potato}{Potato: find the peak crop evapotranspiration, the gross depth per day and the continuous flow per hectare for $E_a=0.75$ and $ET_0=4.0$ mm\,d$^{-1}$.}\sol $ET_c=1.15\times4.0=4.60$ mm\,d$^{-1}$ \ans Gross $=4.60/0.75=6.13$ mm\,d$^{-1}$ \ans Flow $=10\times6.13/24\,\mathrm{m^3\,h^{-1}\,ha^{-1}}=2.56$ m$^3$\,h$^{-1}$, i.e.\ $0.71$ L\,s$^{-1}$\,ha$^{-1}$ \ans

\exbox{Crop example 6b: Soybean}{Soybean: find the peak crop evapotranspiration, the gross depth per day and the continuous flow per hectare for $E_a=0.75$ and $ET_0=4.0$ mm\,d$^{-1}$.}\sol $ET_c=1.15\times4.0=4.60$ mm\,d$^{-1}$ \ans Gross $=4.60/0.75=6.13$ mm\,d$^{-1}$ \ans Flow $=10\times6.13/24\,\mathrm{m^3\,h^{-1}\,ha^{-1}}=2.56$ m$^3$\,h$^{-1}$, i.e.\ $0.71$ L\,s$^{-1}$\,ha$^{-1}$ \ans

\exbox{Crop example 6b: Sugar beet}{Sugar beet: find the peak crop evapotranspiration, the gross depth per day and the continuous flow per hectare for $E_a=0.75$ and $ET_0=4.0$ mm\,d$^{-1}$.}\sol $ET_c=1.20\times4.0=4.80$ mm\,d$^{-1}$ \ans Gross $=4.80/0.75=6.40$ mm\,d$^{-1}$ \ans Flow $=10\times6.40/24\,\mathrm{m^3\,h^{-1}\,ha^{-1}}=2.67$ m$^3$\,h$^{-1}$, i.e.\ $0.74$ L\,s$^{-1}$\,ha$^{-1}$ \ans

\exbox{Crop example 6b: Sunflower}{Sunflower: find the peak crop evapotranspiration, the gross depth per day and the continuous flow per hectare for $E_a=0.75$ and $ET_0=4.0$ mm\,d$^{-1}$.}\sol $ET_c=1.10\times4.0=4.40$ mm\,d$^{-1}$ \ans Gross $=4.40/0.75=5.87$ mm\,d$^{-1}$ \ans Flow $=10\times5.87/24\,\mathrm{m^3\,h^{-1}\,ha^{-1}}=2.44$ m$^3$\,h$^{-1}$, i.e.\ $0.68$ L\,s$^{-1}$\,ha$^{-1}$ \ans

\exbox{Crop example 6b: Onion}{Onion: find the peak crop evapotranspiration, the gross depth per day and the continuous flow per hectare for $E_a=0.75$ and $ET_0=4.0$ mm\,d$^{-1}$.}\sol $ET_c=1.05\times4.0=4.20$ mm\,d$^{-1}$ \ans Gross $=4.20/0.75=5.60$ mm\,d$^{-1}$ \ans Flow $=10\times5.60/24\,\mathrm{m^3\,h^{-1}\,ha^{-1}}=2.33$ m$^3$\,h$^{-1}$, i.e.\ $0.65$ L\,s$^{-1}$\,ha$^{-1}$ \ans

\exbox{Crop example 6b: Alfalfa}{Alfalfa: find the peak crop evapotranspiration, the gross depth per day and the continuous flow per hectare for $E_a=0.75$ and $ET_0=4.0$ mm\,d$^{-1}$.}\sol $ET_c=1.15\times4.0=4.60$ mm\,d$^{-1}$ \ans Gross $=4.60/0.75=6.13$ mm\,d$^{-1}$ \ans Flow $=10\times6.13/24\,\mathrm{m^3\,h^{-1}\,ha^{-1}}=2.56$ m$^3$\,h$^{-1}$, i.e.\ $0.71$ L\,s$^{-1}$\,ha$^{-1}$ \ans

}{}
\IfFileExists{extra/ex06.tex}{\backhead{EXERCISES}{khred}
\begin{exlist}
\item A pipe of 150 mm carries 30 L\,s$^{-1}$. Find the velocity.\exans{1.70 m s$^{-1}$}
\item Hazen--Williams head loss ($C=140$) for $L=400$ m, $D=150$ mm, $Q=30$ L\,s$^{-1}$.\exans{7.1 m}
\item A pump delivers 40 L\,s$^{-1}$ at 55 m head with $\eta=0.7$. Find the power.\exans{30.8 kW}
\item The application efficiency is 0.72 and the conveyance efficiency 0.9. Find the overall efficiency and the gross water for a net need of 400 mm.\exans{0.648; 617 mm}
\item Catch cans: 8, 10, 11, 12, 12, 13, 14, 16 mm. Find CU.\exans{0.854}
\item Emitter flows (L\,h$^{-1}$): 3.6, 3.8, 3.9, 4.0, 4.0, 4.1, 4.2, 4.4. Find EU (lowest value over mean).\exans{0.900}
\item Estimate $R_s$ for $T_{max}=26$, $T_{min}=12$, $R_a=30$, $k_{Rs}=0.16$.\exans{18.0 MJ m$^{-2}$ d$^{-1}$}
\item NDVI for $\rho_{red}=0.06$, $\rho_{NIR}=0.50$.\exans{0.786}
\item A 12 m$^3$\,h$^{-1}$ flow goes through a 63 mm pipe. Is the velocity acceptable?\exans{1.07 m/s; yes}
\item How many cubic metres does 25 mm over 8 ha need?\exans{2000 m$^3$}
\end{exlist}
}{}
\IfFileExists{extra/dyk06.tex}{\dykband{DO YOU KNOW ?}
\begin{exlist}
\item Describe basin, border and furrow irrigation. Where is each used?
\item Distinguish conveyance, application and storage efficiency.
\item What do CU, $DU_{lq}$ and EU measure?
\item State the Hazen--Williams equation and the limits of velocity in a pipe.
\item Why is the optimum pipe diameter found by a cost balance?
\item What checks should be made on weather data before use?
\item How can radiation be estimated from temperature?
\item Explain NDVI and how it gives the cover of a crop.
\end{exlist}
}{}
\IfFileExists{extra/obj06.tex}{\dykband{OBJECTIVE TYPE QUESTIONS}
\begin{exlist}
\item Furrow irrigation is a\opt{surface method}{sprinkler method}{micro-irrigation method}{drainage method}
\item Application efficiency is the share of applied water that is\opt{lost}{stored in the root zone}{evaporated}{pumped}
\item The Hazen--Williams coefficient of new PVC is about\opt{40}{90}{140}{400}
\item Pump power is proportional to\opt{$Q H/\eta$}{$Q/H$}{$\eta/H$}{$QH\eta$}
\item The suitable velocity in a pipe lies between\opt{0 and 0.1}{0.3 and 1.5 m/s}{3 and 10 m/s}{10 and 20 m/s}
\item NDVI of water is\opt{negative}{about 0.9}{about 5}{exactly zero always}
\item Hargreaves' radiation formula uses\opt{rainfall}{temperature range}{wind}{humidity only}
\item Drip systems have efficiencies of about\opt{20 \%}{50 \%}{90 \%}{150 \%}
\item The overall field efficiency is\opt{$E_c+E_a$}{$E_cE_a$}{$E_c/E_a$}{$E_a^2$}
\item Dependable rainfall means\opt{the mean}{rain exceeded with a stated probability}{the maximum}{zero rain}
\end{exlist}
\objkey{1.~(a)\quad 2.~(b)\quad 3.~(c)\quad 4.~(a)\quad 5.~(b)\quad 6.~(a)\quad 7.~(b)\quad 8.~(c)\quad 9.~(b)\quad 10.~(b)}
}{}
\IfFileExists{extra/more06.tex}{\backhead{FURTHER EXERCISES}{khred}
\begin{exlist}
\item A main of 300 m and 150 mm diameter carries 20 L\,s$^{-1}$ (Hazen--Williams $C=130$). (a) Find the friction loss. (b) The outlet is 4 m higher than the pump and needs 25 m of pressure: find the pump head (add 10 per cent for local losses).\exans{(a) 2.86 m; (b) 32.1 m}
\item The pump of the previous exercise has an efficiency of 0.70. Find the power. How much energy does it use to apply 50 mm to 5 ha?\exans{9.0 kW; 313 kWh}
\item A Kostiakov--Lewis intake $Z=k\tau^a+f_0\tau$ has $k=0.008$ m\,min$^{-a}$, $a=0.4$ and $f_0=0.0001$ m\,min$^{-1}$. Find the depth infiltrated in 60 and 240 min.\exans{47.1 mm; 95.6 mm}
\item Manning: a furrow carries 1.5 L\,s$^{-1}$ on a slope of 0.003 with $n=0.04$. For a trapezoid with bottom 0.1 m and side slope 1, the flow depth is 0.08 m. Find the area, the wetted perimeter, the velocity by Manning and the discharge it gives.\exans{0.0144 m$^2$; 0.326 m; 0.171 m\,s$^{-1}$; 2.5 L\,s$^{-1}$}
\item A field 300 m by 200 m receives a gross depth of 70 mm in a set of 12 h. Find the flow needed. How many sets of 1 ha can a supply of 40 L\,s$^{-1}$ irrigate in 7 days of 24 h, at 70 mm each?\exans{97.2 L\,s$^{-1}$; 34 sets}
\end{exlist}
\backhead{ESSAY QUESTIONS}{khred}
{\small\itshape Answer each question in about one page, with equations, a sketch and at least one number from this chapter.}\par\vspace{3pt}
\begin{exlist}
\item Explain the energy equation for flow in pipes, and how the friction and local losses set the size of a pump.
\item Compare the Kostiakov, Kostiakov--Lewis and Green--Ampt descriptions of infiltration. What does each represent physically, and when is each appropriate?
\item Discuss how weather data are put on the twenty-year grid of SWAFL, and the errors that can arise when a short or incomplete record is used.
\item Explain what satellite vegetation indices can and cannot tell about the water status of a field.
\item Discuss the choice of an irrigation system for a field: which properties of the soil, the water source, the crop and the farm decide between surface, drip and sprinkler?
\item Why are uniformity and efficiency different measures? Show with an example how a uniform system can still be inefficient.
\end{exlist}
\dykband{EXERCISES IN SWAFL}
{\small Open SWAFL (desktop or online) with the benchmark project. The chapter of the \href{https://github.com/halilunaziru73-creator/Coupled-Water-Fluxes-Pore-Scale-Dynamics-Precision-Irrigation-Mulching-Feedback-Loops/raw/main/Software/v1.1.0/SWAFL_Desktop/docs/manual_v1.1.0/SWAFL_Manual_v1.1.0.pdf}{SWAFL manual} that explains each task is given in brackets; record your values in a table and compare them with what is expected (\emph{Check}).}\par\vspace{3pt}
\begin{exlist}
\item {\color{khblue}[Manual, Chapter 20]} Find your farm on the satellite map, demarcate it and read the vegetation and water indices. Comment on the variation within the field.\par\noindent{\small\itshape Check: the indices vary with soil and crop cover.}
\item {\color{khblue}[Manual, Chapter 20]} Add a water source and let the program propose an irrigation system for the field. Record the proposal and the reasons it gives.\par\noindent{\small\itshape Check: the proposal depends on the slope, the soil and the source.}
\item {\color{khblue}[Manual, Chapter 19]} Import a weather file of your own and switch back to the benchmark record. Compare the yearly rain and the reference evapotranspiration.\par\noindent{\small\itshape Check: the comparison shows how your climate differs.}
\item {\color{khblue}[Manual, Chapter 18]} Enter the field size as a length and a width and as an area, and check that the design layers give the same volumes.\par\noindent{\small\itshape Check: the volumes agree.}
\end{exlist}
}{}
\swaflvideo{06}
